\documentclass{ieeeaccess}
%DIF LATEXDIFF DIFFERENCE FILE
%DIF DEL oldpaper.tex       Mon Jul 13 19:30:16 2026
%DIF ADD updatedpaper.tex   Fri Jul 17 18:57:16 2026
\usepackage{cite}
\usepackage{amsmath,amssymb,amsfonts}
\usepackage{algorithmic}
\usepackage{graphicx}
\usepackage{textcomp}
\usepackage{caption}
\usepackage{mathptmx}
\usepackage{hyperref}
\usepackage[english]{babel}
\usepackage{microtype} 
\usepackage{xurl}
\usepackage{booktabs}
\usepackage{tabularx}
\usepackage{placeins}
\usepackage{bm}
\usepackage{bbm}
\usepackage{tabularx}
\usepackage{ragged2e}
\makeatletter
\AtBeginDocument{\DeclareMathVersion{bold}
\SetSymbolFont{operators}{bold}{T1}{times}{b}{n}
\SetSymbolFont{NewLetters}{bold}{T1}{times}{b}{it}
\SetMathAlphabet{\mathrm}{bold}{T1}{times}{b}{n}
\SetMathAlphabet{\mathit}{bold}{T1}{times}{b}{it}
\SetMathAlphabet{\mathbf}{bold}{T1}{times}{b}{n}
\SetMathAlphabet{\mathtt}{bold}{OT1}{pcr}{b}{n}
\SetSymbolFont{symbols}{bold}{OMS}{cmsy}{b}{n}
\renewcommand\boldmath{\@nomath\boldmath\mathversion{bold}}}
\makeatother

\def\BibTeX{{\rm B\kern-.05em{\sc i\kern-.025em b}\kern-.08em
    T\kern-.1667em\lower.7ex\hbox{E}\kern-.125emX}}

%Your document starts from here ___________________________________________________
%DIF PREAMBLE EXTENSION ADDED BY LATEXDIFF
%DIF UNDERLINE PREAMBLE %DIF PREAMBLE
\RequirePackage[normalem]{ulem} %DIF PREAMBLE
\RequirePackage{color}\definecolor{RED}{rgb}{1,0,0}\definecolor{BLUE}{rgb}{0,0,1} %DIF PREAMBLE
\providecommand{\DIFaddtex}[1]{{\protect\color{blue}\uwave{#1}}} %DIF PREAMBLE
\providecommand{\DIFdeltex}[1]{{\protect\color{red}\sout{#1}}}                      %DIF PREAMBLE
%DIF SAFE PREAMBLE %DIF PREAMBLE
\providecommand{\DIFaddbegin}{} %DIF PREAMBLE
\providecommand{\DIFaddend}{} %DIF PREAMBLE
\providecommand{\DIFdelbegin}{} %DIF PREAMBLE
\providecommand{\DIFdelend}{} %DIF PREAMBLE
\providecommand{\DIFmodbegin}{} %DIF PREAMBLE
\providecommand{\DIFmodend}{} %DIF PREAMBLE
%DIF FLOATSAFE PREAMBLE %DIF PREAMBLE
\providecommand{\DIFaddFL}[1]{\DIFadd{#1}} %DIF PREAMBLE
\providecommand{\DIFdelFL}[1]{\DIFdel{#1}} %DIF PREAMBLE
\providecommand{\DIFaddbeginFL}{} %DIF PREAMBLE
\providecommand{\DIFaddendFL}{} %DIF PREAMBLE
\providecommand{\DIFdelbeginFL}{} %DIF PREAMBLE
\providecommand{\DIFdelendFL}{} %DIF PREAMBLE
%DIF HYPERREF PREAMBLE %DIF PREAMBLE
\providecommand{\DIFadd}[1]{\texorpdfstring{\DIFaddtex{#1}}{#1}} %DIF PREAMBLE
\providecommand{\DIFdel}[1]{\texorpdfstring{\DIFdeltex{#1}}{}} %DIF PREAMBLE
\newcommand{\DIFscaledelfig}{0.5}
%DIF HIGHLIGHTGRAPHICS PREAMBLE %DIF PREAMBLE
\RequirePackage{settobox} %DIF PREAMBLE
\RequirePackage{letltxmacro} %DIF PREAMBLE
\newsavebox{\DIFdelgraphicsbox} %DIF PREAMBLE
\newlength{\DIFdelgraphicswidth} %DIF PREAMBLE
\newlength{\DIFdelgraphicsheight} %DIF PREAMBLE
% store original definition of \includegraphics %DIF PREAMBLE
\LetLtxMacro{\DIFOincludegraphics}{\includegraphics} %DIF PREAMBLE
\newcommand{\DIFaddincludegraphics}[2][]{{\color{blue}\fbox{\DIFOincludegraphics[#1]{#2}}}} %DIF PREAMBLE
\newcommand{\DIFdelincludegraphics}[2][]{% %DIF PREAMBLE
\sbox{\DIFdelgraphicsbox}{\DIFOincludegraphics[#1]{#2}}% %DIF PREAMBLE
\settoboxwidth{\DIFdelgraphicswidth}{\DIFdelgraphicsbox} %DIF PREAMBLE
\settoboxtotalheight{\DIFdelgraphicsheight}{\DIFdelgraphicsbox} %DIF PREAMBLE
\scalebox{\DIFscaledelfig}{% %DIF PREAMBLE
\parbox[b]{\DIFdelgraphicswidth}{\usebox{\DIFdelgraphicsbox}\\[-\baselineskip] \rule{\DIFdelgraphicswidth}{0em}}\llap{\resizebox{\DIFdelgraphicswidth}{\DIFdelgraphicsheight}{% %DIF PREAMBLE
\setlength{\unitlength}{\DIFdelgraphicswidth}% %DIF PREAMBLE
\begin{picture}(1,1)% %DIF PREAMBLE
\thicklines\linethickness{2pt} %DIF PREAMBLE
{\color[rgb]{1,0,0}\put(0,0){\framebox(1,1){}}}% %DIF PREAMBLE
{\color[rgb]{1,0,0}\put(0,0){\line( 1,1){1}}}% %DIF PREAMBLE
{\color[rgb]{1,0,0}\put(0,1){\line(1,-1){1}}}% %DIF PREAMBLE
\end{picture}% %DIF PREAMBLE
}\hspace*{3pt}}} %DIF PREAMBLE
} %DIF PREAMBLE
\LetLtxMacro{\DIFOaddbegin}{\DIFaddbegin} %DIF PREAMBLE
\LetLtxMacro{\DIFOaddend}{\DIFaddend} %DIF PREAMBLE
\LetLtxMacro{\DIFOdelbegin}{\DIFdelbegin} %DIF PREAMBLE
\LetLtxMacro{\DIFOdelend}{\DIFdelend} %DIF PREAMBLE
\DeclareRobustCommand{\DIFaddbegin}{\DIFOaddbegin \let\includegraphics\DIFaddincludegraphics} %DIF PREAMBLE
\DeclareRobustCommand{\DIFaddend}{\DIFOaddend \let\includegraphics\DIFOincludegraphics} %DIF PREAMBLE
\DeclareRobustCommand{\DIFdelbegin}{\DIFOdelbegin \let\includegraphics\DIFdelincludegraphics} %DIF PREAMBLE
\DeclareRobustCommand{\DIFdelend}{\DIFOaddend \let\includegraphics\DIFOincludegraphics} %DIF PREAMBLE
\LetLtxMacro{\DIFOaddbeginFL}{\DIFaddbeginFL} %DIF PREAMBLE
\LetLtxMacro{\DIFOaddendFL}{\DIFaddendFL} %DIF PREAMBLE
\LetLtxMacro{\DIFOdelbeginFL}{\DIFdelbeginFL} %DIF PREAMBLE
\LetLtxMacro{\DIFOdelendFL}{\DIFdelendFL} %DIF PREAMBLE
\DeclareRobustCommand{\DIFaddbeginFL}{\DIFOaddbeginFL \let\includegraphics\DIFaddincludegraphics} %DIF PREAMBLE
\DeclareRobustCommand{\DIFaddendFL}{\DIFOaddendFL \let\includegraphics\DIFOincludegraphics} %DIF PREAMBLE
\DeclareRobustCommand{\DIFdelbeginFL}{\DIFOdelbeginFL \let\includegraphics\DIFdelincludegraphics} %DIF PREAMBLE
\DeclareRobustCommand{\DIFdelendFL}{\DIFOaddendFL \let\includegraphics\DIFOincludegraphics} %DIF PREAMBLE
%DIF LISTINGS PREAMBLE %DIF PREAMBLE
\RequirePackage{listings} %DIF PREAMBLE
\RequirePackage{color} %DIF PREAMBLE
\lstdefinelanguage{DIFcode}{ %DIF PREAMBLE
%DIF DIFCODE_UNDERLINE %DIF PREAMBLE
  moredelim=[il][\color{red}\sout]{\%DIF\ <\ }, %DIF PREAMBLE
  moredelim=[il][\color{blue}\uwave]{\%DIF\ >\ } %DIF PREAMBLE
} %DIF PREAMBLE
\lstdefinestyle{DIFverbatimstyle}{ %DIF PREAMBLE
	language=DIFcode, %DIF PREAMBLE
	basicstyle=\ttfamily, %DIF PREAMBLE
	columns=fullflexible, %DIF PREAMBLE
	keepspaces=true %DIF PREAMBLE
} %DIF PREAMBLE
\lstnewenvironment{DIFverbatim}{\lstset{style=DIFverbatimstyle}}{} %DIF PREAMBLE
\lstnewenvironment{DIFverbatim*}{\lstset{style=DIFverbatimstyle,showspaces=true}}{} %DIF PREAMBLE
%DIF END PREAMBLE EXTENSION ADDED BY LATEXDIFF

\begin{document}
\history{Date of publication xxxx 00, 0000, date of current version xxxx 00, 0000.}
\doi{10.1109/ACCESS.2024.0429000}


\title{On Reliable Detection of AI-Generated Multimedia Content}
\author{\uppercase{Ansuman Nayak}\authorrefmark{1}, and 
\uppercase{Vijay K. Madisetti}\authorrefmark{2}, \IEEEmembership{Fellow, IEEE}}

\address[1]{College of Computing, Georgia Institute of Technology, Atlanta, GA 30332 USA (e-mail: anayak306@gatech.edu)}
\address[2]{School of Cybersecurity and Privacy, Georgia Institute of Technology, Atlanta, GA 30332 USA (e-mail: vkm@gatech.edu)}

\markboth
{Nayak \headeretal: On Reliable Detection of AI-Generated Multimedia Content}
{Nayak \headeretal: On Reliable Detection of AI-Generated Multimedia Content}

\corresp{Corresponding author: Dr. Vijay Madisetti (e-mail: vkm@gatech.edu).}


% \tfootnote{This work was supported in part by [Funding Source, if applicable].}


\begin{abstract}
With the rapid proliferation of AI-generated multimedia content, the growing need for reliable detection of AI-generated media, including deepfakes and misinformation, has been largely unmet. Uniformly applying large vision-language models (VLMs) to all content can be prohibitively expensive. In this work, we formulate detection as a model selection problem over a set of complementary detectors. We show that disagreement between artifact-based and semantic detectors provides a strong, parameter-free indicator of input difficulty, and that this signal remains stable even when the semantic probe faces deployment shift to unseen generator families. We construct a selective-computation cascade in which lightweight detectors process all inputs and a stronger verifier is invoked only for disagreement cases. On in-domain benchmarks, a well-trained semantic detector achieves the best cost-accuracy tradeoff without VLM escalation. The cascade's utility emerges under generator evolution. Experiments on SDXL and FLUX show that escalating to GPT-5.5 only on disagreement cases recovers up to 4.9 percentage points of accuracy over a stale semantic probe while avoiding 58--65\% of verifier invocations. Cross-model disagreement is a practical routing signal for adaptive computation under distribution shift, requiring no calibrated confidence scores, learned routing policies, or detector retraining.

\end{abstract}

\begin{keywords}
AI-generated image detection, deepfake 
detection, multi-model disagreement, distribution shift, 
vision-language models, forensic routing, CLIP embeddings
\end{keywords}

\titlepgskip=-21pt
\maketitle

% Sections
\section{Introduction}

The rapid advancement of generative models for images and videos has significantly increased the realism and accessibility of synthetic media. Diffusion models~\cite{Wang_2023_ICCV}, GAN-based generators~\cite{wang2020cnn}, and multimodal large language models (MLLMs) can now generate highly convincing content with minimal effort. Although this opens up useful applications, it also raises concerns about misinformation, impersonation, and fraud. As a result, reliably detecting AI-generated media has become an important problem that still needs to be addressed.

One of the main challenges is generalization. Detection models often perform well on the generator families they were trained on, but accuracy degrades significantly when applied to new or unseen generators~\cite{NEURIPS2023_f4d4a021}. Artifact-based detectors exploit low-level frequency and synthesis patterns introduced during generation, while semantic models capture higher-level structural representations. These two signal types fail in different ways, suggesting that combining them provides more robust coverage than either approach alone.


In practical deployment settings, compute cost is a significant concern. High-capacity
vision-language models (VLMs) provide strong reasoning capabilities but are expensive to apply at scale~\cite{chen2024frugalgpt}. Applying these models uniformly to every input is often infeasible in real-time or large-scale systems. Existing approaches typically apply a
fixed detector or ensemble without adapting computation to input difficulty. The central challenge is identifying which inputs require expensive verification rather than applying a
uniform computation to all inputs.


In this work, we address this challenge by modeling AI-generated media detection as a model selection problem over complementary detectors. We find that the disagreement between an artifact-based detector and a semantic detector serves as a proxy for input difficulty and requires no calibration. When the detectors agree, the predictions are highly reliable, whereas disagreements consistently identify cases where cheap detectors are insufficient.

Unlike prior approaches that rely on calibrated confidence scores or learned routing policies, disagreement provides a parameter-free signal that remains reliable under distribution shift. We design a cascade that routes inputs based a signal where lightweight detectors process all inputs, and a stronger verifier is invoked only on disagreement cases. When the semantic probe is trained in-domain, it already achieves strong performance without VLM escalation. The cascade's primary value is under deployment shift, where it recovers
accuracy over a stale semantic probe while avoiding the majority of verifier invocations.This reframes the detection problem as identifying which inputs require expensive computation, rather than uniformly increasing model complexity.

\subsection{Research Questions}

This work is guided by the following research questions:

\begin{itemize}
    \item \textbf{RQ1:} Can we reliably estimate which inputs require more powerful models
    for AI-generated media detection, including under generator distribution shift?
    \item \textbf{RQ2:} Can such estimates be used to guide model selection across heterogeneous detectors without relying on calibrated confidence scores?
    \item \textbf{RQ3:} Does disagreement-based routing reduce verifier cost while recovering accuracy under deployment shift to unseen generator families?
\end{itemize}

\subsection{Contributions}

This work makes the following contributions:

\begin{itemize}
    \item \textbf{C1:} We show that the disagreement between artifact-based and semantic detectors reliably predicts failure cases across generator families and remains
    predictive under deployment shift to unseen generators.

    \item \textbf{C2:} We propose a parameter-free routing strategy that uses disagreement as a difficulty signal for selective model invocation, requiring no confidence
    calibration or policy training.

    \item \textbf{C3:} On deployment shift benchmarks (SDXL and FLUX), cascade escalation to GPT-5.5 recovers up to 4.9 percentage points of accuracy over a stale semantic probe while reducing verifier invocations by 58--65\%.

    \item \textbf{C4:} We show that a learned bandit router does not outperform the disagreement rule, indicating that disagreement captures the dominant routing signal and more complex learned policies offer no meaningful gain.
\end{itemize}


\section{Related Work}
This section reviews prior work related to AI-generated media detection and adaptive inference. We first identify gaps in existing approaches that motivate this work, then discuss detection methods under distribution shift, examine approaches for selective computation, and finally compare routing strategies based on confidence and learned policies.




\subsection{Gaps in Existing Approaches}

Existing forensic detection and cost-aware inference literature leaves several gaps that motivate this work.

\textbf{Single signal reliance.} Most forensic detectors rely on a single signal type, such as diffusion-based artifacts or reconstruction inconsistencies~\cite{Corvi_2023_ICASSP} or semantic embeddings from large pretrained models~\cite{cozzolino2023raising}. These approaches work well within specific generator families but degrade under generator distribution shift. Recent large-scale evaluations show that no single detector consistently generalizes across modern image generation pipelines~\cite{epstein2023online, cozzolino2023raising}.

\textbf{Confidence-based routing under shift.} Cost-aware inference methods and selective classification approaches rely on calibrated confidence estimates to decide when to defer or escalate predictions. Predictive uncertainty in deep models degrades significantly under distribution shift~\cite{ovadia2019can}. As a result, confidence-based routing signals become unreliable in the very regimes where robust detection is most needed. Existing literature on uncertainty estimation and selective prediction does not address this failure mode in heterogeneous detector cascades.

\textbf{Disagreement as a routing signal.} Inter-model disagreement has been studied extensively in active learning~\cite{seung1992query} and ensemble-based uncertainty estimation~\cite{lakshminarayanan2017simple}. In these settings, disagreement identifies informative samples or improves predictive uncertainty. Prior work treats disagreement as a mechanism for aggregation or data selection. Its role as a routing signal for cost-aware inference remains unexplored in forensic detection pipelines. Specifically, using disagreement as a diagnostic to indicate when low-cost detectors should not be trusted has not been studied.

This work addresses these gaps in four ways. First, we combine heterogeneous artifact-based and semantic signals through a parameter-free disagreement rule. Second, we demonstrate empirically that disagreement provides a more reliable routing signal than confidence under distribution shift. We evaluate this at matched operating points across multiple generator families. Third, we account for both latency and token cost when evaluating VLM-based escalation strategies. Fourth, we reinterpret disagreement as a diagnostic for selective invocation of more expensive models, not as a voting mechanism.

\subsection{Detection under Distribution Shift}

A large body of work has focused on detecting AI-generated images using artifact-based features. Early approaches such as CNN-based detectors~\cite{wang2020cnn} exploit low-level patterns introduced by generative models. More recent methods aim to improve generalization across generators, such as universal fake detectors~\cite{ojha2023fakedetect}. Benchmark datasets such as GenImage~\cite{NEURIPS2023_f4d4a021} highlight the difficulty of this problem, showing that many detectors degrade significantly when evaluated on unseen generators.

Other approaches leverage higher-level representations. Models such as CLIP~\cite{pmlr-v139-radford21a} learn semantic features from large-scale image-text data and have been used in downstream detection tasks. Diffusion-specific detection methods, such as DIRE~\cite{Wang_2023_ICCV}, further illustrate the diversity of signals used in modern detection systems. Despite these advances, no single model consistently generalizes across generator families, motivating the use of multiple complementary detectors.

Frequency-domain analysis provides a complementary artifact signal. Spectral methods
exploit periodic and structured artifacts introduced by GAN upsampling and diffusion
denoising pipelines, and have shown strong in-domain generalization~\cite{Corvi_2023_ICASSP}.
Beyond CLIP, recent vision encoders such as DINOv2~\cite{oquab2023dinov2} and
SigLIP2~\cite{tschannen2025siglip2} offer stronger visual representations and have
demonstrated improved robustness under limited training data. Their behavior under
generator distribution shift in forensic detection tasks remains underexplored, motivating
their inclusion as semantic detector candidates in this work.


\subsection{Selective Computation and Model Cascades}

Selective computation has been studied as a way to reduce inference cost by adapting computation to input difficulty. Early work on selective classification~\cite{NIPS2017_4a8423d5} proposes abstaining on uncertain predictions to improve reliability. More recent systems extend this idea to model selection, routing inputs across models of varying capacity. For example, FrugalGPT~\cite{chen2024frugalgpt} reduces cost by selecting among large language models based on estimated utility. These approaches typically rely on confidence scores or calibrated probabilities to guide routing decisions. However, such signals can be unreliable under distribution shift, where models may produce overconfident incorrect predictions.

\subsection{Routing Signals for Model Selection}

Most existing routing strategies use confidence-based signals, such as entropy or prediction probability, to determine whether additional computation is required. While effective in controlled settings, these signals often degrade under distribution shift due to poor calibration. Learned routing policies, such as bandit-based approaches, attempt to improve decision-making by training a model to predict the best action
from input features.

In contrast, this work explores disagreement between heterogeneous detectors as a routing signal. Rather than relying on calibrated confidence scores or learned policies, we use prediction disagreement between artifact-based and semantic models to estimate input difficulty. This provides a simple, calibration-free signal that remains effective under distribution shift.

\begin{table*}[h]
\centering
\small % Reduces font size to help fit more content comfortably
\caption{Comparative analysis of related AI-generated media detection approaches. \checkmark~indicates capability is present; \texttimes~indicates absent.}
\label{tab:litsurvey}
\setlength{\tabcolsep}{4pt} % Slightly tighter spacing between columns

% Define a custom column type 'Y' that is a ragged-right version of 'X'
\newcolumntype{Y}{>{\RaggedRight\arraybackslash}X}

\begin{tabularx}{\textwidth}{p{2.0cm}YYYYYY}
\toprule
\textbf{Feature}
  & \textbf{CNNSpot} \cite{wang2020cnn}
  & \textbf{DIRE} \cite{Wang_2023_ICCV}
  & \textbf{UnivFD} \cite{ojha2023fakedetect}
  & \textbf{BusterX} \cite{wen2025busterxmllmpoweredaigeneratedvideo}
  & \textbf{FrugalGPT} \cite{chen2024frugalgpt}
  & \textbf{Ours} \\
\midrule
Detection signal
  & (\checkmark) Frequency \& upsampling artifacts
  & (\checkmark) Diffusion reconstruction error
  & (\checkmark) CLIP embeddings for generalizable detection
  & (\checkmark) MLLM open-ended visual reasoning
  & (\texttimes) LLM text output, not a forgery detector
  & (\checkmark) Artifact + semantic embeddings + VLM reasoning \\
\midrule
Model architecture
  & (\texttimes) Single CNN/fixed computation
  & (\texttimes) Single diffusion model/fixed computation
  & (\texttimes) Single CLIP backbone/fixed computation
  & (\texttimes) Single MLLM/fixed computation
  & (\checkmark) LLM cascade/variable computation
  & (\checkmark) Heterogeneous three-stage pipeline \\
\midrule
Sequential routing
  & (\texttimes) No routing/uniform inference
  & (\texttimes) No routing/uniform inference
  & (\texttimes) No routing/uniform inference
  & (\texttimes) No routing/uniform inference
  & (\checkmark) Budget-aware cascade routing
  & (\checkmark) Disagreement-based escalation (primary) \\
\midrule
Cost awareness
  & (\texttimes) Fixed cost per sample
  & (\texttimes) Fixed cost per sample
  & (\texttimes) Fixed cost per sample
  & (\texttimes) Fixed cost per sample
  & (\checkmark) Minimizes LLM API cost
  & (\checkmark) Latency and token-aware cascade \\
\midrule
Multi-model disagreement
  & (\texttimes) Single model/no signal
  & (\texttimes) Single model/no signal
  & (\texttimes) Single model/no signal
  & (\texttimes) Single model/no signal
  & (\texttimes) Not applicable to forgery detection
  & (\checkmark) Artifact/semantic detector disagreement as primary signal \\
\midrule
Generator shift eval.
  & (\texttimes) GAN only, degrades on diffusion
  & (\checkmark) Evaluated across diffusion
  & (\checkmark) Evaluated across multiple models
  & (\checkmark) Modern video generators
  & (\texttimes) Not evaluated on media forgery
  & (\checkmark) Video, GAN, diffusion, SDXL, and FLUX generators \\
\midrule
Media type
  & (\texttimes) Images only
  & (\texttimes) Images only
  & (\texttimes) Images only
  & (\checkmark) Video forgery detection
  & (\texttimes) Text only
  & (\checkmark) Video frames and static images \\
\bottomrule
\end{tabularx}
\end{table*}


\section{Proposed Solution}
This section describes the proposed framework for difficulty-aware model selection. We first formalize the routing problem and define the cost model, then introduce the disagreement-based routing strategy, examine the alternative routing baselines used for comparison, and finally detail how latency is measured across routing strategies.

\subsection{Problem Formulation}

We consider the problem of detecting AI-generated media using a set of heterogeneous detectors with different computational costs and failure modes. Given an input sample $x$, the goal is to select an appropriate subset of detectors to apply such that prediction accuracy is maintained while minimizing computational cost.

Let:
\[
\mathcal{M} = \{f_{\text{art}}, f_{\text{sem}}, f_{\text{vlm}}\}
\]
denote the set of available detectors, corresponding to an artifact-based model, a semantic model, and a high-capacity vision-language model. Each model produces a binary prediction $\hat{y}_i \in \{0,1\}$. A routing policy $\pi(x) \in \{0,1\}$ determines whether the VLM is invoked on input $x$, with the artifact and semantic detectors always applied as
low-cost first-stage filters. We evaluate multiple instantiations of this framework in which the artifact detector is instantiated as a dual-stream FFT probe (primary) or CNNSpot
(baseline), and the semantic detector as CLIP, SigLIP2, or DINOv2.

We report on two complementary cost metrics. Wall-clock latency, measured in milliseconds on a fixed GPU, applies to every detection stage:
\[
C_{\text{lat}}(x, \pi) = \kappa^{\text{lat}}_{\text{art}} + \kappa^{\text{lat}}_{\text{sem}}
+ \mathbbm{1}[\pi(x)=1] \cdot \kappa^{\text{lat}}_{\text{vlm}}(x),
\]
where $\kappa^{\text{lat}}_{\text{vlm}}(x)$ is measured per sample because VLM latency depends on input resolution. The token cost, measured as mean VLM input and output tokens per sample, applies only when the VLM is invoked and reflects the standard unit of cost for commercial VLM APIs~\cite{chen2024frugalgpt}:
\[
C_{\text{tok}}(x, \pi) = \mathbbm{1}[\pi(x)=1] \cdot \bigl(\tau^{\text{in}}_{\text{vlm}}(x)
+ \tau^{\text{out}}_{\text{vlm}}(x)\bigr).
\]
The routing objective is to minimize the expected cost under either metric while maintaining a target accuracy $\alpha$:
\[
\min_{\pi} \; \mathbb{E}_x[C(x, \pi)] \quad \text{s.t.} \quad \text{Accuracy}(\pi) \geq \alpha.
\]
We report both $C_{\text{lat}}$ and $C_{\text{tok}}$ in all evaluated methods, as the deployment scenarios vary between self-hosted (latency-bound) and API-based (token-bound) inference.

\subsection{Disagreement as a Difficulty Signal}

We estimate the input difficulty using disagreement between the artifact-based and semantic detectors:

\begin{equation}
d(x) = \mathbbm{1}\left[\hat{y}_{\text{art}}(x) \neq \hat{y}_{\text{sem}}(x)\right].
\DIFaddbegin \label{eq:disagree}
\DIFaddend \end{equation}

When the detectors agree, predictions are highly reliable. Disagreement corresponds to
cases where predictions conflict, indicating that neither signal alone is sufficient and
that additional computation may be required. Unlike confidence-based measures such as
entropy, this signal requires no calibration and shows consistent diagnostic behavior
across generator families. This behavior comes from the complementary failure modes of
both artifact and semantic detectors with their errors are weakly correlated under
distribution shift. This leads to disagreement reliably identifying the cases where cheap detectors
are not trusted.


\subsection{Disagreement-Based Routing}

Using the disagreement signal, we define a simple routing policy:

\begin{equation}
\pi_{\text{disagree}}(x) =
\begin{cases}
0, & \text{if } d(x) = 0 \\
1, & \text{if } d(x) = 1
\end{cases}
\DIFaddbegin \label{eq:pidisagree}
\DIFaddend \end{equation}

where $\pi(x) = 1$ indicates that the VLM is invoked. Inputs where the detectors agree are handled by the cheap stage alone while disagreement cases are escalated to the VLM. This enables selective allocation of computation based on estimated difficulty. Note that outputting the
semantic detector's prediction on all inputs is equivalent to this routing policy without VLM escalation, since the artifact and semantic detectors produce identical predictions on
agreement cases by definition. In-domain, this semantic-only baseline achieves the best
cost-accuracy tradeoff. The cascade's value is under deployment shift, where a stale
semantic probe loses accuracy and VLM escalation on disagreement cases recovers it.

\DIFaddbegin \subsection{\DIFadd{Generalized K-Stage Formulation and Three Operating Regimes}}
\label{sec:kstage}

\DIFadd{The three-detector pipeline evaluated in this paper is an instance of a more
general $K$-stage cascade. Let a sequence of detectors $f_1, \ldots, f_K$
have non-decreasing cost $C_1 \le C_2 \le \cdots \le C_K$ and accuracies
$A_1, \ldots, A_K$ on a fixed evaluation distribution. Define the accuracy
gap between adjacent stages as $\Delta_{i,i+1} = A_{i+1} - A_i$. A cascade
only benefits from escalating stage $i$ to stage $i{+}1$ on a subset of
inputs if $\Delta_{i,i+1}$ restricted to that subset is positive; the
unconditional, dataset-wide comparison of $A_i$ and $A_{i+1}$ is not
sufficient, since escalation is never applied uniformly.
}

\DIFadd{This paper instantiates $K=3$ with $f_1 = f_{\text{art}}$, $f_2 = f_{\text{sem}}$,
and $f_3 = f_{\text{vlm}}$, restricting escalation to the disagreement set
$D = \{x : \hat y_{\text{art}}(x) \neq \hat y_{\text{sem}}(x)\}$
(Eq.~}\eqref{eq:disagree}\DIFadd{). Cascade accuracy decomposes as
}\begin{equation}
\DIFadd{A_{\text{cascade}} = (1-p_D)\, A_{\text{agree}} + p_D\, A_{3\mid D},
\label{eq:decomp}
}\end{equation}
\DIFadd{where $p_D = P(x \in D)$, $A_{\text{agree}}$ is accuracy on the agreement set
(identical for $f_1$ and $f_2$ by construction, Eq.~}\eqref{eq:pidisagree}\DIFadd{),
and $A_{3\mid D}$ is verifier accuracy restricted to $D$. Because the
agreement term is shared with the semantic-only baseline, the sign of
cascade gain over semantic-only depends only on the disagreement subset:
}\begin{equation}
\DIFadd{\text{Gain} = A_{\text{cascade}} - A_2 \;\approx\; p_D \left( A_{3\mid D} - A_{2\mid D} \right),
\label{eq:gain}
}\end{equation}
\DIFadd{where $A_{2\mid D}$ is the semantic detector's own accuracy restricted to
$D$ (typically far below its dataset-wide accuracy, since $D$ is precisely
where $f_2$ is contested). Equation~}\eqref{eq:gain} \DIFadd{formalizes three
qualitative regimes for the two-cheap-detector-plus-verifier setting used
throughout this paper:
}

\textbf{\DIFadd{Case 1 (both cheap stages strong, $A_1 \approx A_2 \approx A_3$).}}
\DIFadd{The artifact and semantic detectors rarely disagree, so $p_D \to 0$ and
Gain $\to 0$ regardless of $A_{3\mid D}$. The verifier's own quality is
nearly irrelevant to end-to-end accuracy because it is invoked on almost
nothing.
}

\textbf{\DIFadd{Case 2 (asymmetric cheap stages, $A_1 \ll A_2$, and $A_2 \gtrsim A_3$).}}
\DIFadd{The disagreement set is dominated by cases where the artifact detector is
wrong and the semantic detector is right. If $A_{3\mid D} < A_{2\mid D}$ --
i.e., the verifier is not reliably better than the semantic detector on
exactly the inputs it disagrees with the artifact detector on -- then
Gain $< 0$: the cascade trades away some of the semantic detector's correct
predictions for the verifier's mixed accuracy on $D$. Any positive delta
over the }\emph{\DIFadd{artifact}}\DIFadd{-only baseline in this regime is attributable to the
semantic detector's quality gap over the artifact detector, not to the
verifier or to the routing mechanism itself.
}

\textbf{\DIFadd{Case 3 (both cheap stages degraded, $A_1 \ll A_3$ and $A_2 \ll A_3$).}}
\DIFadd{Disagreement is common because neither cheap detector is individually
reliable, and if the verifier retains high accuracy on $D$, $A_{3\mid D}$
substantially exceeds $A_{2\mid D}$. This is the only regime in which
Equation~}\eqref{eq:gain} \DIFadd{is reliably positive and large, and the resulting
cascade gain is attributable specifically to escalation rather than to a
quality gap between the two cheap detectors.
}

\textbf{\DIFadd{Breakdown condition.}} \DIFadd{Equation~}\eqref{eq:gain} \DIFadd{also identifies when
the disagreement signal itself degenerates. Disagreement is an informative
routing signal only when (a) the errors of $f_1$ and $f_2$ are weakly
correlated, so that $p_D$ remains large enough to route a meaningful
fraction of the hard cases, and (b) $A_{2\mid D} < A_{3\mid D}$ specifically
on the routed subset. If $f_1$ and $f_2$ share a correlated blind spot --
both systematically wrong on the same inputs -- then $D$ collapses toward
the small residual where the detectors happen to differ, $p_D \to 0$, and
Equation~}\eqref{eq:decomp} \DIFadd{reduces to $A_{\text{cascade}} \approx A_{\text{agree}}$
even though both detectors may be jointly unreliable on the (uncaught)
correlated-error inputs. In this regime, disagreement provides no signal at
all about those inputs; failures caused by correlated blind spots are
invisible to the routing rule by construction, independent of $p_D$
elsewhere in the distribution. Section~VI further connects Cases 1--3 and
this breakdown condition to the empirical results in Section~V.
}

\DIFaddend \subsection{Baseline Routing Strategies}

To evaluate the effectiveness of disagreement-based routing, we compare two alternative strategies.

\textbf{Entropy-based routing.}  
We use the entropy of the artifact detector as a measure of uncertainty:

\begin{equation}
H(x) = - \sum_{i} p_i(x) \log p_i(x),
\end{equation}

where $p_i(x)$ are the predicted class probabilities. Inputs with entropy above a threshold are escalated to the VLM. This represents a standard confidence-based routing approach.

\textbf{Bandit-based routing.}
We train a contextual bandit model using a logistic regression policy on features derived from detector output, including artifact detector confidence, semantic confidence, and prediction disagreement. The model predicts whether an input should be escalated to the VLM, providing a learned routing baseline.

\subsection{Latency Measurement}

Inference latency for each model component is measured independently using synchronized CUDA execution on a fixed hardware setup. The total latency per sample is the sum of the costs of all models invoked by the routing policy. Artifact and semantic detectors are modeled as additive costs, which is a  conservative estimate since these stages could run in parallel in practice. All latency measurements are intended for relative comparison across routing strategies rather than as absolute deployment guarantees.

\FloatBarrier
\begin{figure*}[h!]
    \centering
    \includegraphics[width=0.8\textwidth]{workflowfinal.png}
    \caption{Comparison of detection workflows. (a) FrugalGPT introduces cost-aware cascade
routing for homogeneous LLM outputs. (b) Our framework adapts this paradigm to
heterogeneous forensic detectors, using artifact-semantic disagreement as the primary
routing signal.}
    \label{fig:workflowdiagram}
\end{figure*}
\FloatBarrier

\section{Experiments}

This section describes the evaluation methodology used to assess the proposed routing framework. We first define the compared methods and evaluation metrics, then summarize implementation details and latency measurements.

\subsection{Evaluation Strategy}
We evaluate our framework against artifact-only and semantic-only detectors applied without routing, an always-escalate pipeline that invokes the verifier on every sample, and entropy- and bandit-based routing baselines. UnivFD~\cite{ojha2023fakedetect} is included
as a published 2023 baseline using CLIP ViT-L/14 features trained on GAN images. We
evaluate two verifiers: Qwen2.5-VL-7B~\cite{bai2025qwen25vltechnicalreport} as the
baseline and GPT-5.5~\cite{openai2026gpt55} as the primary, to test whether a stronger verifier better exploits the disagreement signal.
Evaluation focuses on detection accuracy, balanced accuracy, 
average computational cost, and escalation rate. For video 
datasets, a representative frame is sampled from each clip 
and treated as an image detection task following standard 
evaluation practice~\cite{wen2025busterxmllmpoweredaigeneratedvideo}.



\subsection{Implementation}

The pipeline integrates three heterogeneous detection components. The first stage is an artifact-based detector. The primary artifact detector is a dual-stream frequency-domain probe that combines radial FFT magnitude profiles~\cite{dzanic2020fourier} with DCT high-frequency statistics~\cite{frank2020leveraging, Corvi_2023_ICASSP} and color channel
coherence features, trained as a logistic regression classifier. CNNSpot~\cite{wang2020cnn} is retained as an artifact detector baseline. The second stage is a semantic detector. Visual embeddings are extracted with one of three backbones:CLIP
ViT-B/32~\cite{pmlr-v139-radford21a}, SigLIP2~\cite{tschannen2025siglip2}, or
DINOv2~\cite{oquab2023dinov2}. It's then passed through a logistic regression classifier trained on the first 50\% of each dataset. All three semantic probes use the same training split and evaluation protocol. The third and final stage is an optional VLM verifier which is only invoked 
on disagreement cases.


To utilize GPT-5.5 as the primary image verifier, we prompt it as follows:

\noindent
\begin{minipage}{\linewidth}
\vspace{10pt}
\itshape
``Analyze this image carefully. Determine whether it is a REAL photograph or an AI-generated FAKE image. Consider both possibilities equally. Answer with exactly one word: REAL or FAKE.''
\vspace{10pt}
\end{minipage}

When applied to the real images bundled with AIGIBench (PASCAL VOC 2007--2012
photographs), GPT-5.5 classified 82--97\% of them as \emph{FAKE} despite being
genuine photographs. PASCAL VOC images are archival scans with heavy JPEG
compression that a frontier VLM trained on modern high-quality imagery does not
recognize as real, reducing real-class recall to near zero. We replaced these
images with ImageNet val/nature photographs from the GenImage SD~1.4
partition~\cite{NEURIPS2023_f4d4a021}, which GPT-5.5 correctly identifies as real
and which match the real-image source used to train the SD~1.4 semantic probes.
Any VLM-based evaluation on AIGIBench should substitute the bundled real images
with contemporary photographs to avoid this calibration artifact. \DIFaddbegin \DIFadd{This
substitution also means the real half of the SDXL and FLUX test sets is
in-distribution for the semantic probes by construction; Section~VI.G
quantifies the effect this has on the reported deployment-shift recovery.
}\DIFaddend 


Confidence scores are derived from the logistic probe output for the FFT artifact detector and all semantic detectors. For the CNNSpot baseline, confidence is the maximum softmax probability from the ResNet output. The bandit and entropy baselines use these scores as features. The disagreement rule would only use binary predictions and would require no confidence
calibration.

The routing policy operates on the joint output of the artifact and semantic detectors. 
In the disagreement threshold configuration, escalation is triggered whenever the two first-stage detectors produce conflicting 
predictions. In the entropy threshold configuration, escalation is 
triggered when the artifact detectors prediction entropy exceeds $\tau = 0.5$. The 
bandit router applies a logistic regression policy trained on artifact
confidence, semantic confidence, and their disagreement signal to produce 
a binary escalation decision.

We measure deployment latency on an NVIDIA RTX A4500 GPU with synchronized CUDA execution. The FFT probe takes 2.5~ms per sample. Each semantic detector takes 10.4~ms per sample,
giving a cheap-stage latency of 12.9~ms. This represents a 4.4$\times$ reduction in cheap-stage cost compared to the CNNSpot-based configuration (56.2~ms). GPT-5.5 latency on the original benchmarks averages 1{,}595~ms on GenBuster, 1{,}459~ms on SD~1.4, and 1{,}572~ms on BigGAN. On the SDXL and FLUX benchmarks, where images are higher resolution, GPT-5.5 latency averages 3{,}084--3{,}103~ms and 3{,}670--3{,}849~ms respectively. Qwen2.5-VL
latency on the original benchmarks is 190~ms (GenBuster), 83~ms (SD~1.4), and 71~ms
(BigGAN) under always-escalate. We report both latency and token cost throughout our results. Latency comparisons assume sequential FFT and semantic detector execution. Parallel
execution would further reduce cheap-stage cost.

\section{Results}

This section presents empirical results across video, diffusion, and GAN benchmarks. We first describe the datasets and evaluation protocol, then compare detector performance, analyze routing behavior, and finally study the robustness of the proposed signal through ablations.

\subsection{Experimental Setup}

We evaluate on three in-domain benchmarks spanning distinct generator families.
GenBuster-200K-mini~\cite{wen2025busterxmllmpoweredaigeneratedvideo} comprises 2{,}000
samples from eight modern AI video generators and 1{,}000 real video clips.
GenImage SD~1.4~\cite{NEURIPS2023_f4d4a021} comprises 6{,}000 Stable Diffusion v1.4
images and 6{,}000 real ImageNet images. GenImage
BigGAN~\cite{NEURIPS2023_f4d4a021} comprises 6{,}000 BigGAN images and 6{,}000 real
ImageNet images. Each sample is represented by a single frame or image following standard
practice~\cite{wen2025busterxmllmpoweredaigeneratedvideo}.

We additionally evaluate under deployment shift on two modern diffusion benchmarks.
Fake images are drawn from AIGIBench~\cite{li2025artificial} SDXL and FLUX1-dev
subsets. Real images are ImageNet val/nature photographs from the GenImage SD~1.4
partition. AIGIBench's bundled PASCAL VOC real images are not used because they
cause systematic VLM calibration failure (Section~IV.B). Each benchmark contains
4{,}000 samples (2{,}000 fake, 2{,}000 real). All semantic probes are trained
exclusively on SD~1.4 data. No SDXL or FLUX images appear in any training split.


For the in-domain benchmarks, each dataset is partitioned into three non-overlapping splits
using stratified sampling. The first split (50\%) trains the semantic probe classifiers,
with a held-out 20\% subset used for probe validation. The remaining 50\% forms the routing
policy training and evaluation pool, from which a held-out 40\% subset is used for all
reported evaluations. No detector or routing policy is evaluated on samples seen during
training. Evaluation sets contain 400 samples for GenBuster and 480 samples each for
SD~1.4 and BigGAN.

We compare FFT-only, CNNSpot-only~\cite{wang2020cnn}, UnivFD~\cite{ojha2023fakedetect},
and semantic-only configurations using CLIP~\cite{pmlr-v139-radford21a},
SigLIP2~\cite{tschannen2025siglip2}, and DINOv2~\cite{oquab2023dinov2}. Routing
baselines include entropy threshold ($\tau = 0.5$), disagreement-based escalation to Qwen2.5-VL-7B or GPT-5.5, and uniform escalation to Qwen2.5-VL. All configurations
are evaluated under identical conditions on the same held-out splits.

All accuracy figures are reported with 95\% bootstrap confidence intervals over 1{,}000
resamples. \DIFaddbegin \DIFadd{Deltas between paired configurations (cascade vs.\ stale probe, cascade vs.\
FFT-only) are additionally assessed with a paired bootstrap over 10}{\DIFadd{,}}\DIFadd{000 resamples,
which resamples indices once per draw and recomputes both arms on the same resample so
correlation between paired errors is preserved (Section~V.D.6). }\DIFaddend Balanced accuracy is the mean of fake and real recall. Latency is measured in
milliseconds on an NVIDIA RTX A4500 GPU.

\subsection{Detector Performance}


\textbf{Artifact detectors.} Table~\DIFdelbegin \DIFdel{\ref{tab:artifact_comparison} }\DIFdelend \DIFaddbegin \DIFadd{\ref{tab:main_results} }\DIFaddend compares the dual-stream
FFT probe against CNNSpot across all three in-domain benchmarks. FFT achieves 75.3\%,
74.8\%, and 99.8\% accuracy on GenBuster, SD~1.4, and BigGAN respectively, against
CNNSpot's 52.0\%, 50.4\%, and 72.1\%. This gain comes at 2.5~ms per sample versus
CNNSpot's 45.8~ms, an 18$\times$ reduction in artifact-stage latency. The FFT probe
maintains balanced recall across fake and real classes, while CNNSpot exhibits strong real
bias on out-of-domain generators (4.3\% fake recall on GenBuster).

\textbf{Semantic detectors.} Among the three semantic probes, SigLIP2 achieves the
strongest in-domain accuracy on all benchmarks: 98.5\% on GenBuster, 92.7\% on SD~1.4,
and 99.6\% on BigGAN. CLIP and DINOv2 are competitive (96.3--99.4\% on BigGAN, 91.3--96.5\% on GenBuster), with DINOv2 trailing on SD~1.4 (89.0\%). 
All three probes are trained on in-domain data and run at 10.4~ms per sample.

When semantic probes are trained in-domain, the semantic-only baseline achieves the best
cost-accuracy tradeoff on all three benchmarks. Routing configurations that invoke
Qwen2.5-VL on disagreement cases achieve 89--90.5\% on GenBuster and 84--87\% on SD~1.4,
below the semantic-only baselines (96.3--98.5\% and 89--92.7\% respectively). We expected this since
Qwen provides weak signal on these benchmarks (69.0\%, 60.6\%, 70.4\% under
uniform escalation), and invoking it on the hardest cases cannot recover to the level of a
well-trained in-domain probe. \DIFaddbegin \DIFadd{This is the expected signature of }\textbf{\DIFadd{Case 2}} \DIFadd{in the
K-stage formulation (Section~III.D): the artifact detector is comparatively weak
($A_1 = 75$--$99.8\%$) while the semantic detector is comparatively strong
($A_2 = 89$--$98.5\%$) and is not exceeded by the verifier on the disagreement subset
specifically. Restricted to the FFT+SigLIP2 disagreement subset, SigLIP2's own accuracy
$A_{2\mid D}$ is 95.1\% on GenBuster and 82.6\% on SD~1.4, both well above Qwen's accuracy
on that same subset ($A_{3\mid D}=$ 63.6\% and 60.6\% respectively, Table~\ref{tab:verifier_comparison}).
%DIF >  NOTE: A_{2|D} sourced from routing/results/e0_fft_disagreement_stats.csv
%DIF >  (dis_sem_acc, SigLIP2 rows, pre-existing data). A_{3|D} matches
%DIF >  Table~\ref{tab:verifier_comparison}'s Qwen column, both now fully fixed
%DIF >  this session (routing/fix_table6_verifier_comparison.py +
%DIF >  routing/fill_table6_gpt55_gap.py).
The cascade's positive delta over FFT-only in this regime
(Table~\ref{tab:crossdataset}) reflects the semantic detector's quality gap over the
artifact detector, not a genuine contribution from VLM escalation; the value
proposition of escalation holds only under Case 3 (Section~VI).
}\DIFaddend 

\textbf{Verifier comparison.} Table~\ref{tab:verifier_comparison} compares Qwen2.5-VL
and GPT-5.5 on the \DIFdelbegin \DIFdel{disagreement subset. }\DIFdelend \DIFaddbegin \DIFadd{FFT+SigLIP2 disagreement subset -- the same subset the primary
cascade escalates throughout this paper, rather than a subset defined by a different
detector pair. On this subset, Qwen2.5-VL is }\emph{\DIFadd{more}} \DIFadd{accurate than }\DIFaddend GPT-5.5 \DIFdelbegin \DIFdel{achieves 91.9\% accuracy on
the
GenBuster disagreement subset versus Qwen's 75.0}\DIFdelend \DIFaddbegin \DIFadd{on
GenBuster (63.6\% vs.\ 55.5}\DIFaddend \%, a \DIFdelbegin \DIFdel{16.9 percentage point gain}\DIFdelend \DIFaddbegin \DIFadd{$-$8.1 point GPT-5.5 delta): the GenBuster disagreement
cases that FFT and SigLIP2 disagree on turn out to be ones where GPT-5.5 does
comparatively poorly, the opposite of the ordering suggested by a subset built from a
weaker artifact detector}\DIFaddend . On SD~1.4, GPT-5.5\DIFdelbegin \DIFdel{is modestly better (63.9\% vs58.8\%)}\DIFdelend \DIFaddbegin \DIFadd{'s advantage over Qwen widens substantially
on the correct subset (77.3\% vs}\DIFaddend .\DIFdelbegin \DIFdel{On BigGAN}\DIFdelend \DIFaddbegin \DIFadd{\ 60.6\%}\DIFaddend , \DIFaddbegin \DIFadd{$+$16.7 points, versus a $+$5.0-point gap on
the old subset). BigGAN's disagreement subset is only 3 samples (FFT and SigLIP2 agree
on 99.4\% of BigGAN inputs, Table~\ref{tab:cheap_path}), too few to support any claim
about relative verifier accuracy there; both verifiers happen to score 66.7\% on this
subset. This dataset-dependent reversal -- }\DIFaddend GPT-5.5 \DIFdelbegin \DIFdel{performs below
Qwen (43.0\% vs 71.1\% ). We attribute this
reversal to the same distribution mismatch
documented in Section~IV.B. The images in BigGAN have a distinct visual style consistent with early GAN
synthesis that the latest VLMs do not associate with AI generation. GPT-5.5's strong
performance on modern generators (GenBuster video, SDXL, FLUX) confirms it is better
calibrated to current generation artifacts}\DIFdelend \DIFaddbegin \DIFadd{is not uniformly stronger than Qwen,
even restricted to the disagreement cases the cascade actually escalates -- is itself
evidence that verifier choice should be validated per-deployment rather than assumed
from a single benchmark, and reinforces that Case~2/Case~3 identity depends on which
verifier is paired with which detector pair on which dataset}\DIFaddend .

\textbf{UnivFD baseline.} UnivFD~\cite{ojha2023fakedetect}, trained on GAN images using
CLIP ViT-L/14, achieves zero fake recall on GenBuster and near-zero on SD~1.4 (5.9\%),
matching chance accuracy. On BigGAN, where the test distribution matches its training
domain, it achieves 91.5\%. GAN-trained detectors fail categorically on video generators
and diffusion images regardless of backbone strength. The advantage of trainable semantic
probes over fixed pretrained ones is consistent across all evaluated backbones.



\begin{figure*}[t]
    \centering
    \includegraphics[width=\textwidth]{plot_accuracy_comparison.png}
    \caption{Detector accuracy across three benchmarks. Semantic probes dominate in-domain on all datasets. FFT substantially outperforms CNNSpot on diffusion and video generators. The cascade (FFT + SigLIP2 + Qwen) falls below semantic-only in-domain, confirming that routing adds value only when the semantic probe faces deployment shift.}

    \label{fig:accuracy_comparison}
\end{figure*}


%DIF <  --- Table: Artifact Detector Comparison (E2) ---
\DIFdelbegin %DIFDELCMD < \begin{table}[h]
%DIFDELCMD < \large
%DIFDELCMD < \centering
%DIFDELCMD < %%%
%DIFDELCMD < \caption{%
{%DIFAUXCMD
\DIFdelFL{Artifact detector comparison. FFT outperforms CNNSpot on all benchmarks at
18$\times$ lower latency.}}
%DIFAUXCMD
%DIFDELCMD < \label{tab:artifact_comparison}
%DIFDELCMD < \resizebox{\columnwidth}{!}{%
%DIFDELCMD < \begin{tabular}{lcccccc}
%DIFDELCMD < \toprule
%DIFDELCMD < Detector & GenBuster & SD~1.4 & BigGAN & Fake Rec. (GB) & Latency (ms) \\
%DIFDELCMD < \midrule
%DIFDELCMD < CNNSpot~\cite{wang2020cnn} & 0.520 & 0.502 & 0.721 & 0.043 & 45.8 \\
%DIFDELCMD < FFT (ours)                  & 0.753 & 0.748 & 0.998 & 0.677 & 2.5  \\
%DIFDELCMD < \bottomrule
%DIFDELCMD < \end{tabular}%
%DIFDELCMD < }
%DIFDELCMD < \end{table}
%DIFDELCMD < 

%DIFDELCMD < %%%
%DIF <  --- Table: GenBuster main results ---
%DIFDELCMD < \begin{table}[h]
%DIFDELCMD < \large
%DIFDELCMD < \centering
%DIFDELCMD < %%%
%DIFDELCMD < \caption{%
{%DIFAUXCMD
\DIFdelFL{Detection results on GenBuster-200K-mini. CI denotes 95\% bootstrap confidence
interval.}}
%DIFAUXCMD
%DIFDELCMD < \label{tab:genbuster}
%DIFDELCMD < \resizebox{\columnwidth}{!}{%
%DIFDELCMD < \begin{tabular}{lcccc}
%DIFDELCMD < \toprule
%DIFDELCMD < Method & Accuracy & 95\% CI & Latency (ms) & Esc. Rate \\
%DIFDELCMD < \midrule
%DIFDELCMD < FFT-only          & 0.753 & [0.710, 0.793] &  2.5  & 0.00 \\
%DIFDELCMD < CNNSpot-only      & 0.520 & [0.470, 0.568] & 45.8  & 0.00 \\
%DIFDELCMD < UnivFD            & 0.512 & [0.463, 0.565] & 56.2  & 0.00 \\
%DIFDELCMD < \midrule
%DIFDELCMD < CLIP-only         & 0.965 & [0.945, 0.983] & 10.4  & 0.00 \\
%DIFDELCMD < SigLIP2-only      & 0.985 & [0.970, 0.995] & 10.4  & 0.00 \\
%DIFDELCMD < DINOv2-only       & 0.963 & [0.945, 0.978] & 10.4  & 0.00 \\
%DIFDELCMD < \midrule
%DIFDELCMD < Entropy thresh    & 0.527 & [0.480, 0.570] & 17.9  & 0.03 \\
%DIFDELCMD < Disagree→Qwen (FFT+CLIP)    & 0.900 & [0.870, 0.928] & 55.1  & 0.24 \\
%DIFDELCMD < Disagree→Qwen (FFT+SigLIP2) & 0.905 & [0.875, 0.933] & 57.9  & 0.26 \\
%DIFDELCMD < Disagree→Qwen (FFT+DINOv2)  & 0.890 & [0.860, 0.918] & 58.4  & 0.26 \\
%DIFDELCMD < Always escalate   & 0.690 & [0.648, 0.738] & 190.2 & 1.00 \\
%DIFDELCMD < \bottomrule
%DIFDELCMD < \end{tabular}%
%DIFDELCMD < }
%DIFDELCMD < \end{table}
%DIFDELCMD < 

%DIFDELCMD < %%%
%DIF <  --- Table: SD 1.4 main results ---
%DIFDELCMD < \begin{table}[h]
%DIFDELCMD < \large
%DIFDELCMD < \centering
%DIFDELCMD < %%%
%DIFDELCMD < \caption{%
{%DIFAUXCMD
\DIFdelFL{Detection results on GenImage SD~1.4.}}
%DIFAUXCMD
%DIFDELCMD < \label{tab:sd14}
%DIFDELCMD < \resizebox{\columnwidth}{!}{%
%DIFDELCMD < \begin{tabular}{lcccc}
%DIFDELCMD < \toprule
%DIFDELCMD < Method & Accuracy & 95\% CI & Latency (ms) & Esc. Rate \\
%DIFDELCMD < \midrule
%DIFDELCMD < FFT-only          & 0.748 & [0.710, 0.790] &  2.5  & 0.00 \\
%DIFDELCMD < CNNSpot-only      & 0.504 & [0.460, 0.546] & 45.8  & 0.00 \\
%DIFDELCMD < UnivFD            & 0.529 & [0.485, 0.575] & 56.2  & 0.00 \\
%DIFDELCMD < \midrule
%DIFDELCMD < CLIP-only         & 0.913 & [0.888, 0.935] & 10.4  & 0.00 \\
%DIFDELCMD < SigLIP2-only      & 0.927 & [0.904, 0.950] & 10.4  & 0.00 \\
%DIFDELCMD < DINOv2-only       & 0.890 & [0.860, 0.919] & 10.4  & 0.00 \\
%DIFDELCMD < \midrule
%DIFDELCMD < Entropy thresh    & 0.525 & [0.479, 0.571] & 14.3  & 0.07 \\
%DIFDELCMD < Disagree→Qwen (FFT+CLIP)    & 0.865 & [0.833, 0.896] & 36.9  & 0.30 \\
%DIFDELCMD < Disagree→Qwen (FFT+SigLIP2) & 0.867 & [0.835, 0.898] & 31.9  & 0.28 \\
%DIFDELCMD < Disagree→Qwen (FFT+DINOv2)  & 0.842 & [0.808, 0.873] & 32.5  & 0.29 \\
%DIFDELCMD < Always escalate   & 0.606 & [0.565, 0.654] & 83.2  & 1.00 \\
%DIFDELCMD < \bottomrule
%DIFDELCMD < \end{tabular}%
%DIFDELCMD < }
%DIFDELCMD < \end{table}
%DIFDELCMD < 

%DIFDELCMD < %%%
%DIF <  --- Table: BigGAN main results ---
%DIFDELCMD < \begin{table}[h]
%DIFDELCMD < %%%
\DIFdelendFL %DIF >  --- Table: Consolidated main results (was 4 separate tables: artifact_comparison,
%DIF >  genbuster, sd14, biggan -- merged this session to address reviewer point (7),
%DIF >  table sprawl. Artifact-detector-only comparison (CNNSpot vs FFT) is now a
%DIF >  subset of rows within this table rather than a separate table. ---
\DIFaddbeginFL \begin{table*}[t]
\DIFaddendFL \large
\centering
\caption{Detection results \DIFaddbeginFL \DIFaddFL{across all three in-domain benchmarks. CI denotes 95\%
bootstrap confidence interval. FFT outperforms CNNSpot }\DIFaddendFL on \DIFdelbeginFL \DIFdelFL{GenImage BigGAN}\DIFdelendFL \DIFaddbeginFL \DIFaddFL{every benchmark at
18$\times$ lower latency (rows 1-2 of each block)}\DIFaddendFL . \DIFaddbeginFL \DIFaddFL{BigGAN: }\DIFaddendFL FFT and semantic
detectors near-saturate \DIFdelbeginFL \DIFdelFL{.
}\DIFdelendFL \DIFaddbeginFL \DIFaddFL{and }\DIFaddendFL VLM escalation is rarely triggered ($<$1.3\%).}
\DIFdelbeginFL %DIFDELCMD < \label{tab:biggan}
%DIFDELCMD < \resizebox{\columnwidth}{!}{%
%DIFDELCMD < \begin{tabular}{lcccc}
%DIFDELCMD < \toprule
%DIFDELCMD < Method & Accuracy & 95\% CI & Latency (ms) & Esc. Rate \\
%DIFDELCMD < \midrule
%DIFDELCMD < FFT-only          & 0.998 & [0.994, 1.000] &  2.5  & 0.00 \\
%DIFDELCMD < CNNSpot-only      & 0.721 & [0.681, 0.758] & 45.8  & 0.00 \\
%DIFDELCMD < UnivFD            & 0.915 & [0.890, 0.938] & 56.2  & 0.00 \\
%DIFDELCMD < \midrule
%DIFDELCMD < CLIP-only         & 0.994 & [0.985, 1.000] & 10.4  & 0.00 \\
%DIFDELCMD < SigLIP2-only      & 0.996 & [0.990, 1.000] & 10.4  & 0.00 \\
%DIFDELCMD < DINOv2-only       & 0.990 & [0.979, 0.998] & 10.4  & 0.00 \\
%DIFDELCMD < \midrule
%DIFDELCMD < Entropy thresh    & 0.733 & [0.692, 0.771] & 12.9  & 0.06 \\
%DIFDELCMD < Disagree→Qwen (FFT+CLIP)    & 0.998 & [0.994, 1.000] & 13.3  & 0.01 \\
%DIFDELCMD < Disagree→Qwen (FFT+SigLIP2) & 0.998 & [0.994, 1.000] & 13.2  & 0.01 \\
%DIFDELCMD < Disagree→Qwen (FFT+DINOv2)  & 0.998 & [0.994, 1.000] & 13.5  & 0.01 \\
%DIFDELCMD < Always escalate   & 0.704 & [0.665, 0.745] & 70.6  & 1.00 \\
%DIFDELCMD < \bottomrule
%DIFDELCMD < \end{tabular}%
%DIFDELCMD < }
%DIFDELCMD < \end{table}
%DIFDELCMD < %%%
\DIFdelend \DIFaddbegin \label{tab:main_results}
\begin{tabular}{llcccc}
\toprule
\DIFadd{Dataset }& \DIFadd{Method }& \DIFadd{Accuracy }& \DIFadd{95\% CI }& \DIFadd{Latency (ms) }& \DIFadd{Esc. Rate }\\
\midrule
\DIFadd{GenBuster }& \DIFadd{FFT-only          }& \DIFadd{0.753 }& [\DIFadd{0.710, 0.793}] &  \DIFadd{2.5  }& \DIFadd{0.00 }\\
\DIFadd{GenBuster }& \DIFadd{CNNSpot-only      }& \DIFadd{0.520 }& [\DIFadd{0.470, 0.568}] & \DIFadd{45.8  }& \DIFadd{0.00 }\\
\DIFadd{GenBuster }& \DIFadd{UnivFD            }& \DIFadd{0.512 }& [\DIFadd{0.463, 0.565}] & \DIFadd{56.2  }& \DIFadd{0.00 }\\
\DIFadd{GenBuster }& \DIFadd{CLIP-only         }& \DIFadd{0.965 }& [\DIFadd{0.945, 0.983}] & \DIFadd{10.4  }& \DIFadd{0.00 }\\
\DIFadd{GenBuster }& \DIFadd{SigLIP2-only      }& \DIFadd{0.985 }& [\DIFadd{0.970, 0.995}] & \DIFadd{10.4  }& \DIFadd{0.00 }\\
\DIFadd{GenBuster }& \DIFadd{DINOv2-only       }& \DIFadd{0.963 }& [\DIFadd{0.945, 0.978}] & \DIFadd{10.4  }& \DIFadd{0.00 }\\
\DIFadd{GenBuster }& \DIFadd{Entropy thresh    }& \DIFadd{0.527 }& [\DIFadd{0.480, 0.570}] & \DIFadd{17.9  }& \DIFadd{0.03 }\\
\DIFadd{GenBuster }& \DIFadd{Disagree→Qwen (FFT+CLIP)    }& \DIFadd{0.900 }& [\DIFadd{0.870, 0.928}] & \DIFadd{55.1  }& \DIFadd{0.24 }\\
\DIFadd{GenBuster }& \DIFadd{Disagree→Qwen (FFT+SigLIP2) }& \DIFadd{0.905 }& [\DIFadd{0.875, 0.933}] & \DIFadd{57.9  }& \DIFadd{0.26 }\\
\DIFadd{GenBuster }& \DIFadd{Disagree→Qwen (FFT+DINOv2)  }& \DIFadd{0.890 }& [\DIFadd{0.860, 0.918}] & \DIFadd{58.4  }& \DIFadd{0.26 }\\
\DIFadd{GenBuster }& \DIFadd{Always escalate   }& \DIFadd{0.690 }& [\DIFadd{0.648, 0.738}] & \DIFadd{190.2 }& \DIFadd{1.00 }\\
\midrule
\DIFadd{SD~1.4 }& \DIFadd{FFT-only          }& \DIFadd{0.748 }& [\DIFadd{0.710, 0.790}] &  \DIFadd{2.5  }& \DIFadd{0.00 }\\
\DIFadd{SD~1.4 }& \DIFadd{CNNSpot-only      }& \DIFadd{0.504 }& [\DIFadd{0.460, 0.546}] & \DIFadd{45.8  }& \DIFadd{0.00 }\\
\DIFadd{SD~1.4 }& \DIFadd{UnivFD            }& \DIFadd{0.529 }& [\DIFadd{0.485, 0.575}] & \DIFadd{56.2  }& \DIFadd{0.00 }\\
\DIFadd{SD~1.4 }& \DIFadd{CLIP-only         }& \DIFadd{0.913 }& [\DIFadd{0.888, 0.935}] & \DIFadd{10.4  }& \DIFadd{0.00 }\\
\DIFadd{SD~1.4 }& \DIFadd{SigLIP2-only      }& \DIFadd{0.927 }& [\DIFadd{0.904, 0.950}] & \DIFadd{10.4  }& \DIFadd{0.00 }\\
\DIFadd{SD~1.4 }& \DIFadd{DINOv2-only       }& \DIFadd{0.890 }& [\DIFadd{0.860, 0.919}] & \DIFadd{10.4  }& \DIFadd{0.00 }\\
\DIFadd{SD~1.4 }& \DIFadd{Entropy thresh    }& \DIFadd{0.525 }& [\DIFadd{0.479, 0.571}] & \DIFadd{14.3  }& \DIFadd{0.07 }\\
\DIFadd{SD~1.4 }& \DIFadd{Disagree→Qwen (FFT+CLIP)    }& \DIFadd{0.865 }& [\DIFadd{0.833, 0.896}] & \DIFadd{36.9  }& \DIFadd{0.30 }\\
\DIFadd{SD~1.4 }& \DIFadd{Disagree→Qwen (FFT+SigLIP2) }& \DIFadd{0.867 }& [\DIFadd{0.835, 0.898}] & \DIFadd{31.9  }& \DIFadd{0.28 }\\
\DIFadd{SD~1.4 }& \DIFadd{Disagree→Qwen (FFT+DINOv2)  }& \DIFadd{0.842 }& [\DIFadd{0.808, 0.873}] & \DIFadd{32.5  }& \DIFadd{0.29 }\\
\DIFadd{SD~1.4 }& \DIFadd{Always escalate   }& \DIFadd{0.606 }& [\DIFadd{0.565, 0.654}] & \DIFadd{83.2  }& \DIFadd{1.00 }\\
\midrule
\DIFadd{BigGAN }& \DIFadd{FFT-only          }& \DIFadd{0.998 }& [\DIFadd{0.994, 1.000}] &  \DIFadd{2.5  }& \DIFadd{0.00 }\\
\DIFadd{BigGAN }& \DIFadd{CNNSpot-only      }& \DIFadd{0.721 }& [\DIFadd{0.681, 0.758}] & \DIFadd{45.8  }& \DIFadd{0.00 }\\
\DIFadd{BigGAN }& \DIFadd{UnivFD            }& \DIFadd{0.915 }& [\DIFadd{0.890, 0.938}] & \DIFadd{56.2  }& \DIFadd{0.00 }\\
\DIFadd{BigGAN }& \DIFadd{CLIP-only         }& \DIFadd{0.994 }& [\DIFadd{0.985, 1.000}] & \DIFadd{10.4  }& \DIFadd{0.00 }\\
\DIFadd{BigGAN }& \DIFadd{SigLIP2-only      }& \DIFadd{0.996 }& [\DIFadd{0.990, 1.000}] & \DIFadd{10.4  }& \DIFadd{0.00 }\\
\DIFadd{BigGAN }& \DIFadd{DINOv2-only       }& \DIFadd{0.990 }& [\DIFadd{0.979, 0.998}] & \DIFadd{10.4  }& \DIFadd{0.00 }\\
\DIFadd{BigGAN }& \DIFadd{Entropy thresh    }& \DIFadd{0.733 }& [\DIFadd{0.692, 0.771}] & \DIFadd{12.9  }& \DIFadd{0.06 }\\
\DIFadd{BigGAN }& \DIFadd{Disagree→Qwen (FFT+CLIP)    }& \DIFadd{0.998 }& [\DIFadd{0.994, 1.000}] & \DIFadd{13.3  }& \DIFadd{0.01 }\\
\DIFadd{BigGAN }& \DIFadd{Disagree→Qwen (FFT+SigLIP2) }& \DIFadd{0.998 }& [\DIFadd{0.994, 1.000}] & \DIFadd{13.2  }& \DIFadd{0.01 }\\
\DIFadd{BigGAN }& \DIFadd{Disagree→Qwen (FFT+DINOv2)  }& \DIFadd{0.998 }& [\DIFadd{0.994, 1.000}] & \DIFadd{13.5  }& \DIFadd{0.01 }\\
\DIFadd{BigGAN }& \DIFadd{Always escalate   }& \DIFadd{0.704 }& [\DIFadd{0.665, 0.745}] & \DIFadd{70.6  }& \DIFadd{1.00 }\\
\bottomrule
\end{tabular}
\end{table*}
\DIFaddend 

% --- Table: Verifier comparison (E4) ---
\begin{table}[h]
\centering
\small
\setlength{\tabcolsep}{3pt}  % reduced from 4pt to 3pt
\caption{Verifier accuracy on the \DIFaddbeginFL \DIFaddFL{FFT+SigLIP2 }\DIFaddendFL disagreement subset \DIFdelbeginFL \DIFdelFL{. GPT-5.5 substantially outperforms
Qwen on modern generator types but underperforms on BigGAN, consistent with calibration
to modern imagery}\DIFdelendFL \DIFaddbeginFL \DIFaddFL{(matching the
disagreement set used by the primary cascade throughout this paper)}\DIFaddendFL . Since Qwen is
baseline here, improvement is marked as ---\DIFaddbeginFL \DIFaddFL{. $n$ is small on BigGAN because FFT and
SigLIP2 rarely disagree there (Case 1).}\DIFaddendFL }
\label{tab:verifier_comparison}
\begin{tabular}{llcccc}
\toprule
Dataset & Verifier & Acc. & Improvement & $n$ & Lat. (ms) \\
\midrule
GenBuster & Qwen2.5-VL-7B  & \DIFdelbeginFL \DIFdelFL{0.750 }\DIFdelendFL \DIFaddbeginFL \DIFaddFL{0.636 }\DIFaddendFL & ---        & \DIFdelbeginFL \DIFdelFL{172 }\DIFdelendFL \DIFaddbeginFL \DIFaddFL{110 }\DIFaddendFL & \DIFdelbeginFL \DIFdelFL{178 }\DIFdelendFL \DIFaddbeginFL \DIFaddFL{174 }\DIFaddendFL \\
GenBuster & GPT-5.5        & \DIFdelbeginFL \DIFdelFL{0.919 }\DIFdelendFL \DIFaddbeginFL \DIFaddFL{0.555 }\DIFaddendFL & \DIFdelbeginFL \DIFdelFL{$+$0.169   }\DIFdelendFL \DIFaddbeginFL \DIFaddFL{$-$0.081   }\DIFaddendFL & \DIFdelbeginFL \DIFdelFL{172 }\DIFdelendFL \DIFaddbeginFL \DIFaddFL{110 }\DIFaddendFL & \DIFdelbeginFL \DIFdelFL{1595 }\DIFdelendFL \DIFaddbeginFL \DIFaddFL{2443 }\DIFaddendFL \\
\midrule
SD~1.4    & Qwen2.5-VL-7B  & \DIFdelbeginFL \DIFdelFL{0.588 }\DIFdelendFL \DIFaddbeginFL \DIFaddFL{0.606 }\DIFaddendFL & ---        & \DIFdelbeginFL \DIFdelFL{238 }\DIFdelendFL \DIFaddbeginFL \DIFaddFL{132 }\DIFaddendFL & \DIFdelbeginFL \DIFdelFL{74 }\DIFdelendFL \DIFaddbeginFL \DIFaddFL{69 }\DIFaddendFL \\
SD~1.4    & GPT-5.5        & \DIFdelbeginFL \DIFdelFL{0.639 }\DIFdelendFL \DIFaddbeginFL \DIFaddFL{0.773 }\DIFaddendFL & $+$\DIFdelbeginFL \DIFdelFL{0.050   }\DIFdelendFL \DIFaddbeginFL \DIFaddFL{0.167   }\DIFaddendFL & \DIFdelbeginFL \DIFdelFL{238 }\DIFdelendFL \DIFaddbeginFL \DIFaddFL{132 }\DIFaddendFL & \DIFdelbeginFL \DIFdelFL{1459 }\DIFdelendFL \DIFaddbeginFL \DIFaddFL{2492 }\DIFaddendFL \\
\midrule
BigGAN    & Qwen2.5-VL-7B  & \DIFdelbeginFL \DIFdelFL{0.711 }\DIFdelendFL \DIFaddbeginFL \DIFaddFL{0.667 }\DIFaddendFL & ---        &   \DIFdelbeginFL \DIFdelFL{135 }\DIFdelendFL \DIFaddbeginFL \DIFaddFL{3 }\DIFaddendFL & \DIFdelbeginFL \DIFdelFL{58 }\DIFdelendFL \DIFaddbeginFL \DIFaddFL{53 }\DIFaddendFL \\
BigGAN    & GPT-5.5        & \DIFdelbeginFL \DIFdelFL{0.430 }\DIFdelendFL \DIFaddbeginFL \DIFaddFL{0.667 }\DIFaddendFL & \DIFdelbeginFL \DIFdelFL{$-$0.281   }\DIFdelendFL \DIFaddbeginFL \DIFaddFL{$+$0.000   }\DIFaddendFL &   \DIFdelbeginFL \DIFdelFL{135 }\DIFdelendFL \DIFaddbeginFL \DIFaddFL{3 }\DIFaddendFL & \DIFdelbeginFL \DIFdelFL{1572 }\DIFdelendFL \DIFaddbeginFL \DIFaddFL{2023 }\DIFaddendFL \\
\bottomrule
\end{tabular}
\end{table}
%DIF >  NOTE: Table rebuilt on the correct FFT+SigLIP2 disagreement subset (was
%DIF >  CNNSpot+CLIP, n=172/238/135) this session. Source:
%DIF >  routing/results/table6_fft_siglip2_fix.csv (fix_table6_verifier_comparison.py)
%DIF >  + routing/fill_table6_gpt55_gap.py (142 new GPT-5.5 calls to complete
%DIF >  coverage). GenBuster's verifier ordering flips entirely (GPT-5.5 no longer
%DIF >  beats Qwen there); SD~1.4's GPT-5.5 advantage widens; BigGAN's n=3 is too
%DIF >  small to support the old "GPT-5.5 underperforms on BigGAN" claim.


% --- Table: E8 SDXL/FLUX deployment shift main results ---
\DIFdelbegin %DIFDELCMD < \begin{table}[h]
%DIFDELCMD < %%%
\DIFdelendFL \DIFaddbeginFL \begin{table*}[t]
\DIFaddendFL \large
\centering
\caption{Deployment shift results on SDXL and FLUX. Semantic probes are trained on
SD~1.4 and tested without retraining. \DIFaddbeginFL \DIFaddFL{FFT is also trained on SD~1.4 and tested without
retraining (stale, like the semantic probe). }\DIFaddendFL Cascade escalates to \DIFdelbeginFL \DIFdelFL{GPT-5.5 }\DIFdelendFL \DIFaddbeginFL \DIFaddFL{the named verifier }\DIFaddendFL on
disagreement cases. Gain = cascade accuracy $-$ stale semantic accuracy\DIFaddbeginFL \DIFaddFL{.}\DIFaddendFL }
\label{tab:e8_results}
\DIFdelbeginFL %DIFDELCMD < \resizebox{\columnwidth}{!}{%
%DIFDELCMD < \begin{tabular}{llccc}
%DIFDELCMD < \toprule
%DIFDELCMD < Dataset & Semantic & Stale Sem. & Cascade (GPT-5.5) & Gain \\
%DIFDELCMD < \midrule
%DIFDELCMD < SDXL & CLIP    & 0.872 & 0.893 & $+$0.021 \\
%DIFDELCMD < SDXL & SigLIP2 & 0.847 & 0.881 & $+$0.034 \\
%DIFDELCMD < SDXL & DINOv2  & 0.855 & 0.891 & $+$0.036 \\
%DIFDELCMD < \midrule
%DIFDELCMD < FLUX & CLIP    & 0.920 & 0.914 & $-$0.006 \\
%DIFDELCMD < FLUX & SigLIP2 & 0.909 & 0.913 & $+$0.004 \\
%DIFDELCMD < FLUX & DINOv2  & 0.815 & 0.864 & $+$0.049 \\
%DIFDELCMD < \bottomrule
%DIFDELCMD < \end{tabular}%
%DIFDELCMD < }
%DIFDELCMD < \end{table}
%DIFDELCMD < %%%
\DIFdelend \DIFaddbegin \begin{tabular}{llcccccc}
\toprule
\DIFadd{Dataset }& \DIFadd{Semantic }& \DIFadd{Stale FFT }& \DIFadd{Stale Sem. }& \DIFadd{Casc.\ (Qwen) }& \DIFadd{Gain (Qwen) }& \DIFadd{Casc.\ (GPT-5.5) }& \DIFadd{Gain (GPT-5.5) }\\
\midrule
\DIFadd{SDXL }& \DIFadd{CLIP    }& \DIFadd{0.622 }& \DIFadd{0.872 }& \DIFadd{0.826 }& \DIFadd{$-$0.047 }& \DIFadd{0.893 }& \DIFadd{$+$0.021 }\\
\DIFadd{SDXL }& \DIFadd{SigLIP2 }& \DIFadd{0.622 }& \DIFadd{0.847 }& \DIFadd{0.817 }& \DIFadd{$-$0.030 }& \DIFadd{0.881 }& \DIFadd{$+$0.034 }\\
\DIFadd{SDXL }& \DIFadd{DINOv2  }& \DIFadd{0.622 }& \DIFadd{0.855 }& \DIFadd{0.824 }& \DIFadd{$-$0.031 }& \DIFadd{0.891 }& \DIFadd{$+$0.036 }\\
\midrule
\DIFadd{FLUX }& \DIFadd{CLIP    }& \DIFadd{0.578 }& \DIFadd{0.920 }& \DIFadd{0.779 }& \DIFadd{$-$0.140 }& \DIFadd{0.914 }& \DIFadd{$-$0.006 }\\
\DIFadd{FLUX }& \DIFadd{SigLIP2 }& \DIFadd{0.578 }& \DIFadd{0.909 }& \DIFadd{0.779 }& \DIFadd{$-$0.130 }& \DIFadd{0.913 }& \DIFadd{$+$0.004 }\\
\DIFadd{FLUX }& \DIFadd{DINOv2  }& \DIFadd{0.578 }& \DIFadd{0.815 }& \DIFadd{0.734 }& \DIFadd{$-$0.081 }& \DIFadd{0.864 }& \DIFadd{$+$0.049 }\\
\bottomrule
\end{tabular}
\end{table*}
%DIF >  NOTE: Stale FFT column added from routing/results/e8_modern_diffusion.csv
%DIF >  (fft_stale_acc column) -- pre-existing data from before this session, not
%DIF >  a new experiment. Qwen cascade/gain columns added this session from a new
%DIF >  Qwen2.5-VL-7B run on all 3991 unique SDXL/FLUX disagreement-subset images
%DIF >  (routing/run_qwen_sdxl_flux.py, routing/results/e8_qwen_cascade.csv).
%DIF >  Addresses reviewer request to add FFT-only and Qwen2.5-only columns to
%DIF >  Tables 7/8. Note Qwen gives NEGATIVE gain on all six configs here (unlike
%DIF >  GPT-5.5's consistent positive gain) -- this deployment-shift cascade's
%DIF >  Case-3 status depends specifically on using a strong verifier (GPT-5.5),
%DIF >  not on the routing mechanism alone; a weak verifier (Qwen) turns this into
%DIF >  a Case-2-like result even though both cheap detectors are genuinely
%DIF >  degraded. Worth a sentence in Section VI.E (VLM Domain Sensitivity) or
%DIF >  wherever this table is first discussed in Results. Stale FFT (62.2%/57.8%)
%DIF >  is well below GPT-5.5's cascade and stale-semantic accuracy here (unlike
%DIF >  the cross-family Table 18 result) -- with GPT-5.5 as verifier this
%DIF >  deployment-shift setting is a genuine Case 3, not Case 2, consistent with
%DIF >  Section VI.B; with Qwen as verifier it behaves like Case 2 instead.
\DIFaddend 

% --- Table: E8 token savings ---
\begin{table}[h]
\centering
\caption{Disagreement rates and verifier invocation savings on SDXL and FLUX. Token
savings = 1 $-$ escalation rate.}
\label{tab:e8_efficiency}
\begin{tabular}{llcc}
\toprule
Dataset & Semantic & Disagree Rate & Token Savings \\
\midrule
SDXL & CLIP    & 35.2\% & 64.8\% \\
SDXL & SigLIP2 & 36.6\% & 63.4\% \\
SDXL & DINOv2  & 38.6\% & 61.4\% \\
\midrule
FLUX & CLIP    & 40.3\% & 59.7\% \\
FLUX & SigLIP2 & 41.4\% & 58.6\% \\
FLUX & DINOv2  & 41.4\% & 58.6\% \\
\bottomrule
\end{tabular}
\end{table}



\subsection{Routing Analysis}

In this section we analyze whether artifact-semantic disagreement between the FFT artifact detector and
the semantic probe provides a reliable routing signal, and how selective escalation
affects the cost-accuracy tradeoff across in-domain and deployment-shift conditions.

Figure~\ref{fig:agree_disagree} shows disagreement rates across all five evaluation
datasets. On in-domain benchmarks, the rate ranges from 24--30\%. On BigGAN it collapses
to under 1\%, consistent with the FFT probe near-saturating at 99.8\% accuracy and
leaving almost no room for artifact-semantic conflict. Under deployment shift on SDXL and
FLUX, the rate rises to 35--42\%, marking the conditions where the stale semantic probe
loses accuracy and VLM escalation provides gain.


\begin{figure*}[t]
    \centering
    \includegraphics[width=\textwidth]{plot_fig_c_disagreement_rate.png}
    \caption{Artifact-semantic disagreement rate across five datasets. The artifact detector
is the FFT probe. Semantic detectors are CLIP, SigLIP2, and DINOv2. Disagreement stays
at 24--30\% on in-domain benchmarks and collapses to under 1\% on BigGAN where the FFT
probe near-saturates. The hatched SDXL and FLUX bars mark deployment shift, where
disagreement rises sharply and VLM escalation provides the most gain.}
    \label{fig:agree_disagree}
\end{figure*}


Table~\ref{tab:crossdataset} shows cascade gain over the FFT-only baseline across datasets.

\begin{table}[h]
\centering
\caption{Cascade accuracy and gain over FFT-only across datasets\DIFaddbeginFL \DIFaddFL{, alongside the
individual FFT-only, SigLIP2-only, and Qwen (always-escalate) baselines}\DIFaddendFL . Cascade is
FFT + SigLIP2 with Qwen escalation on disagreement cases. \DIFaddbeginFL \DIFaddFL{Complete baseline columns
make explicit that in-domain cascade gain over FFT-only tracks the FFT/SigLIP2
accuracy gap (Case 2, Section~III.D), not Qwen's own accuracy.}\DIFaddendFL }
\label{tab:crossdataset}
\small
\DIFdelbeginFL %DIFDELCMD < \begin{tabularx}{\columnwidth}{l l c c c}
%DIFDELCMD < %%%
\DIFdelendFL \DIFaddbeginFL \begin{tabularx}{\columnwidth}{l l c c c c c}
\DIFaddendFL \toprule
\textbf{Dataset} & \textbf{Generator} & \textbf{FFT} & \DIFaddbeginFL \textbf{\DIFaddFL{SigLIP2}} & \textbf{\DIFaddFL{Qwen}} & \DIFaddendFL \textbf{Cascade} & \textbf{Gain} \\
\midrule
GenBuster & Video     & 0.753 & \DIFaddbeginFL \DIFaddFL{0.985 }& \DIFaddFL{0.690 }& \DIFaddendFL 0.905 & $+$0.152 \\
SD~1.4    & Diffusion & 0.748 & \DIFaddbeginFL \DIFaddFL{0.927 }& \DIFaddFL{0.606 }& \DIFaddendFL 0.867 & $+$0.119 \\
BigGAN    & GAN       & 0.998 & \DIFaddbeginFL \DIFaddFL{0.996 }& \DIFaddFL{0.704 }& \DIFaddendFL 0.998 & $+$0.000 \\
\bottomrule
\end{tabularx}
\end{table}
%DIF >  NOTE: SigLIP2 and Qwen (always-escalate) columns added from
%DIF >  routing/results/e0_fft_table.csv (Semantic-only and Always-Qwen rows,
%DIF >  SigLIP2 backbone) -- pre-existing data, no new experiment needed. Addresses
%DIF >  reviewer request to complete baseline columns in this table.


When the FFT artifact detector and semantic probe agree, the agreement subset achieves
above 96\% accuracy across all three in-domain benchmarks (Table~\ref{tab:cheap_path}).
When they disagree, FFT accuracy on the disagreement subset drops to under 18\% on
GenBuster and SD~1.4, and neither detector is reliable on its own. This asymmetry
confirms that artifact-semantic disagreement is a stable routing signal without requiring
confidence calibration.

Table~\ref{tab:escalation} breaks down escalation behavior by case type on GenBuster for
the FFT+SigLIP2 configuration. The disagreement rule escalates all 103 disagreement cases
and bypasses all 297 agreement cases by construction. FFT is correct on 296 of 297
agreement-subset samples. The single case where both detectors agree but FFT is wrong
represents the fundamental limit of the routing signal. When both cheap detectors fail
simultaneously, no artifact-semantic disagreement can flag the error. At 99.7\% agreement-
subset accuracy, this limit affects under 0.3\% of all inputs. The disagreement rule
identifies 98 of the 99 total FFT failures on GenBuster.

\begin{table}[h]
\centering
\caption{Escalation behavior by case type on GenBuster-200K-mini (FFT+SigLIP2).
The disagreement rule covers 98 of 99 total FFT failures. The one uncatchable failure
is the case where both cheap detectors simultaneously err.}
\label{tab:escalation}
\begin{tabular}{lcc}
\toprule
Case & $n$ & Disagree esc. \\
\midrule
FFT correct, models agree    & 296 & 0.00 \\
FFT correct, models disagree &   5 & 1.00 \\
FFT wrong, models agree      &   1 & 0.00 \\
FFT wrong, models disagree   &  98 & 1.00 \\
\bottomrule
\end{tabular}
\end{table}

Table~\ref{tab:per_generator} breaks down accuracy by generator on GenBuster.
CNNSpot achieves zero fake recall on all eight video generators, confirming that
GAN-trained detectors fail categorically on modern video content. Disagreement-routed
escalation recovers accuracy across generators but falls below the semantic-only
baseline, consistent with Qwen's limited accuracy on video frames.


\begin{table}[h]
\centering
\caption{Per-generator detection accuracy on GenBuster-200K-mini. Disagreement-routed
escalation improves over CNNSpot across all generators but falls below CLIP, reflecting
Qwen's limited accuracy on video frames}
\label{tab:per_generator}
\begin{tabular}{lrcccc}
\toprule
Generator & $n$ & CNNSpot & Cascade & CLIP \\
\midrule
Gen3   & 24  & 0.000 & 0.792 & 1.000 \\
Jimeng & 22  & 0.136 & 0.318 & 0.818 \\
Kling  & 20  & 0.050 & 0.500 & 0.800 \\
Luma   & 18  & 0.222 & 0.667 & 0.778 \\
Pika   & 15  & 0.000 & 0.600 & 0.933 \\
Sora   & 31  & 0.032 & 0.613 & 0.968 \\
Vidu   & 38  & 0.026 & 0.421 & 0.789 \\
Wanx   & 27  & 0.000 & 0.519 & 0.741 \\
\midrule
Real   & 205 & 0.966 & 0.956 & 0.966 \\
\bottomrule
\end{tabular}
\end{table}


The retrained bandit router, trained on the same \DIFdelbegin \DIFdel{Qwen verifier used in }\DIFdelend \DIFaddbegin \DIFadd{FFT+SigLIP2+Qwen features as }\DIFaddend the disagreement threshold configuration, converges to performance that matches the disagreement rule on \DIFdelbegin \DIFdel{both GenBuster (0.755 vs0.755}\DIFdelend \DIFaddbegin \DIFadd{all three datasets: SD~1.4 (0.867 vs.\ 0.867}\DIFaddend ) and BigGAN (\DIFdelbegin \DIFdel{0.925 vs0.925).On SD~1.4 the bandit trails disagreement by 11.9\% (0.623 vs0.742) , reflecting the weaker Qwen signal on that dataset, which provides a noisier reward for policy learning}\DIFdelend \DIFaddbegin \DIFadd{0.998 vs.\ 0.998) exactly, and GenBuster (0.917 vs.\ 0.917) on an internal held-out split of its validation set (Table~\ref{tab:transfer} footnote; GenBuster's Qwen predictions do not cover the training-split frames the standard protocol requires)}\DIFaddend . The bandit's \DIFaddbegin \DIFadd{exact }\DIFaddend convergence to disagreement on \DIFdelbegin \DIFdel{two of three datasets }\DIFdelend \DIFaddbegin \DIFadd{every dataset }\DIFaddend suggests that disagreement captures \DIFdelbegin \DIFdel{most of the useful routing signal for routing in }\DIFdelend the \DIFaddbegin \DIFadd{dominant routing signal in this }\DIFaddend clean-signal regime, and that more complex learned policies offer \DIFdelbegin \DIFdel{little }\DIFdelend \DIFaddbegin \DIFadd{no }\DIFaddend benefit over the parameter-free rule.


\DIFdelbegin %DIFDELCMD < \begin{table}[h]
%DIFDELCMD < \centering
%DIFDELCMD < %%%
%DIFDELCMD < \caption{%
{%DIFAUXCMD
\DIFdelFL{Mean per-sample latency and accuracy on GenBuster-200K-mini. Latency reduction
is relative to always-escalate to Qwen (190.2~ms)}}
%DIFAUXCMD
%DIFDELCMD < \label{tab:latency}
%DIFDELCMD < \begin{tabular}{lccc}
%DIFDELCMD < \toprule
%DIFDELCMD < %%%
\DIFdelFL{Method }%DIFDELCMD < & %%%
\DIFdelFL{Latency (ms) }%DIFDELCMD < & %%%
\DIFdelFL{Reduction }%DIFDELCMD < & %%%
\DIFdelFL{Accuracy }%DIFDELCMD < \\
%DIFDELCMD < \midrule
%DIFDELCMD < %%%
\DIFdelFL{FFT-only               }%DIFDELCMD < &   %%%
\DIFdelFL{2.5 }%DIFDELCMD < & %%%
\DIFdelFL{99\%  }%DIFDELCMD < & %%%
\DIFdelFL{0.753 }%DIFDELCMD < \\
%DIFDELCMD < %%%
\DIFdelFL{SigLIP2-only           }%DIFDELCMD < &  %%%
\DIFdelFL{10.4 }%DIFDELCMD < & %%%
\DIFdelFL{95\%  }%DIFDELCMD < & %%%
\DIFdelFL{0.985 }%DIFDELCMD < \\
%DIFDELCMD < %%%
\DIFdelFL{Entropy thresh         }%DIFDELCMD < &  %%%
\DIFdelFL{17.9 }%DIFDELCMD < & %%%
\DIFdelFL{91\%  }%DIFDELCMD < & %%%
\DIFdelFL{0.527 }%DIFDELCMD < \\
%DIFDELCMD < %%%
\DIFdelFL{Disagree }\DIFdelendFL \DIFaddbeginFL \DIFaddFL{Relative to always-escalate to Qwen (190.2~ms per sample, Table~\ref{tab:main_results}),
the disagreement cascade }\DIFaddendFL (FFT+SigLIP2) \DIFdelbeginFL %DIFDELCMD < &  %%%
\DIFdelFL{57.9 }%DIFDELCMD < & %%%
\DIFdelendFL \DIFaddbeginFL \DIFaddFL{reduces mean per-sample latency by }\DIFaddendFL 70\% \DIFdelbeginFL %DIFDELCMD < & %%%
\DIFdelFL{0.905 }%DIFDELCMD < \\
%DIFDELCMD < %%%
\DIFdelFL{Always escalate        }%DIFDELCMD < & %%%
\DIFdelFL{190.2 }%DIFDELCMD < &  %%%
\DIFdelFL{0\% }%DIFDELCMD < & %%%
\DIFdelFL{0.690 }%DIFDELCMD < \\
%DIFDELCMD < \bottomrule
%DIFDELCMD < \end{tabular}
%DIFDELCMD < \end{table}
%DIFDELCMD < %%%
\DIFdelend \DIFaddbegin \DIFadd{on
GenBuster (57.9~ms), and the semantic-only baseline alone reduces it by 95\% (10.4~ms).
%DIF >  NOTE: Former tab:latency table deleted this session -- fully redundant with
%DIF >  Table~\ref{tab:main_results} rows (same methods/numbers), addressing
%DIF >  reviewer point (7). Only new information (reduction %) folded into prose.
}\DIFaddend 


\begin{table}[h]
\centering
\caption{Disagreement vs.\ entropy routing at matched escalation rates (CNNSpot+CLIP
configuration). Disagreement outperforms entropy by 16--19 percentage points on every
dataset}
\label{tab:matched_entropy}
\begin{tabular}{lcccc}
\toprule
Dataset & Method & Accuracy & Esc. Rate & Gap \\
\midrule
GenBuster & Disagree thresh      & 0.755 & 0.44 & --     \\
GenBuster & Entropy matched      & 0.590 & 0.26 & $-$0.165 \\
\midrule
SD 1.4    & Disagree thresh      & 0.742 & 0.49 & --     \\
SD 1.4    & Entropy matched      & 0.550 & 0.43 & $-$0.192 \\
\midrule
BigGAN    & Disagree thresh      & 0.925 & 0.28 & --     \\
BigGAN    & Entropy matched      & 0.760 & 0.28 & $-$0.165 \\
\bottomrule
\end{tabular}
% \footnotetext{Disagreement threshold values reflect the FFT+SigLIP2+Qwen configuration.
% Entropy-matched comparison uses matched escalation rates from the CNNSpot artifact detector;
% the ordering is expected to hold with FFT entropy given the larger disagreement accuracy gap.}
The FFT-based disagreement rule achieves substantially higher absolute accuracy on
these same datasets (Table~\ref{tab:crossdataset}); the entropy comparison here is
retained to show the signal-quality advantage holds independent of the artifact
detector used.
\end{table}


\begin{figure*}[h]
    \centering
    \includegraphics[width=\textwidth]{plot_cost_accuracy_all.png}
    \caption{Cost--accuracy tradeoff per semantic backbone. In-domain (left), semantic-only probes
are Pareto-optimal and the cascade adds latency without accuracy gain. Under deployment
shift on SDXL (right), stale probe accuracy drops 4--8 percentage points from in-domain
and the cascade with GPT-5.5 partially recovers it through selective escalation.}
    \label{fig:cost_accuracy}
\end{figure*}


\subsection{Ablation Studies}
The ablation studies characterize cheap-path coverage, robustness of the
disagreement-based routing signal, and cascade behavior under cross-family probe
transfer.


\subsubsection{Cheap-Path Coverage}

Table~\ref{tab:cheap_path} quantifies cheap-path coverage for the FFT+SigLIP2
configuration. On every dataset, at least 74\% of samples fall in the agreement subset where FFT accuracy exceeds 96\%, and escalation is limited to the remaining disagreement
cases. On BigGAN, the agreement rate reaches 99.4\% and cheap-path accuracy is 100\%,
leaving only 3 samples for VLM escalation.

\FloatBarrier
\begin{table}[h]
\centering
\caption{Agreement-subset coverage and accuracy for the FFT+SigLIP2 configuration.
Agreement-subset accuracy exceeds 96\% on all three datasets.}

\label{tab:cheap_path}
\begin{tabular}{lcccc}
\toprule
Dataset & Agree & Cheap Acc & Disagree & Esc. \\
        & $n$   &           & $n$      & rate \\
\midrule
GenBuster & 297 & 0.997 & 103 & 0.258 \\
SD~1.4    & 348 & 0.966 & 132 & 0.275 \\
BigGAN    & 477 & 1.000 &   3 & 0.006 \\
\bottomrule
\end{tabular}
\end{table}
\FloatBarrier

\DIFaddbegin \subsubsection{\DIFadd{Probe-Degradation Sweep}}
\label{sec:limited_data}

\DIFadd{To directly validate the Case~2-to-Case~3 transition (Section~VI.B) within a single
experiment rather than relying on cross-dataset comparison, we hold the FFT artifact
detector fixed and sweep semantic-probe quality by varying the fraction of labeled
training data available to the probe (1\%, 5\%, 10\%, 25\%, 50\%, 100\%), evaluating
semantic-only accuracy and FFT+semantic cascade accuracy (Qwen verifier) at each point.
Table~\ref{tab:limited_data} reports SD~1.4 with SigLIP2 and DINOv2 backbones.
}

\begin{table}[h]
\centering
\caption{\DIFaddFL{Probe-degradation sweep on SD~1.4. Cascade gain (Casc.\ $-$ Sem.) changes sign
as the semantic probe weakens, crossing from Case~2 (negative gain) to Case~3 (positive
gain) between 5\% and 10\% training data.}}
\label{tab:limited_data}
\small
\begin{tabular}{lcccc}
\toprule
\DIFaddFL{Backbone }& \DIFaddFL{Train frac. }& \DIFaddFL{Sem.-only }& \DIFaddFL{Cascade }& \DIFaddFL{Gain }\\
\midrule
\DIFaddFL{SigLIP2  }& \DIFaddFL{1\%   }& \DIFaddFL{0.740 }& \DIFaddFL{0.794 }& \DIFaddFL{$+$0.054 }\\
\DIFaddFL{SigLIP2  }& \DIFaddFL{5\%   }& \DIFaddFL{0.842 }& \DIFaddFL{0.833 }& \DIFaddFL{$-$0.008 }\\
\DIFaddFL{SigLIP2  }& \DIFaddFL{10\%  }& \DIFaddFL{0.898 }& \DIFaddFL{0.854 }& \DIFaddFL{$-$0.044 }\\
\DIFaddFL{SigLIP2  }& \DIFaddFL{25\%  }& \DIFaddFL{0.908 }& \DIFaddFL{0.863 }& \DIFaddFL{$-$0.046 }\\
\DIFaddFL{SigLIP2  }& \DIFaddFL{50\%  }& \DIFaddFL{0.919 }& \DIFaddFL{0.869 }& \DIFaddFL{$-$0.050 }\\
\DIFaddFL{SigLIP2  }& \DIFaddFL{100\% }& \DIFaddFL{0.927 }& \DIFaddFL{0.867 }& \DIFaddFL{$-$0.060 }\\
\midrule
\DIFaddFL{DINOv2   }& \DIFaddFL{1\%   }& \DIFaddFL{0.669 }& \DIFaddFL{0.769 }& \DIFaddFL{$+$0.100 }\\
\DIFaddFL{DINOv2   }& \DIFaddFL{5\%   }& \DIFaddFL{0.781 }& \DIFaddFL{0.798 }& \DIFaddFL{$+$0.017 }\\
\DIFaddFL{DINOv2   }& \DIFaddFL{10\%  }& \DIFaddFL{0.808 }& \DIFaddFL{0.810 }& \DIFaddFL{$+$0.002 }\\
\DIFaddFL{DINOv2   }& \DIFaddFL{25\%  }& \DIFaddFL{0.829 }& \DIFaddFL{0.808 }& \DIFaddFL{$-$0.021 }\\
\DIFaddFL{DINOv2   }& \DIFaddFL{50\%  }& \DIFaddFL{0.852 }& \DIFaddFL{0.821 }& \DIFaddFL{$-$0.031 }\\
\DIFaddFL{DINOv2   }& \DIFaddFL{100\% }& \DIFaddFL{0.890 }& \DIFaddFL{0.842 }& \DIFaddFL{$-$0.048 }\\
\bottomrule
\end{tabular}
\end{table}

\DIFadd{Both backbones cross from positive to negative cascade gain as training data increases
from 1\% to 10\%, i.e., as the semantic probe strengthens from well below the FFT
detector's 74.8\% SD~1.4 accuracy to well above it. GenBuster does not exhibit this
crossover within the tested range: even a probe trained on 1\% of GenBuster's data (16
labeled examples) already exceeds FFT's 75.3\% accuracy there, so gain stays negative at
every fraction, illustrating that the crossover point is set by the FFT detector's own
strength on a given generator family, not by data quantity in isolation. BigGAN similarly
does not cross over, staying in Case~1 (small positive gain, near-zero disagreement rate)
throughout, since FFT already exceeds 99\% accuracy independent of semantic-probe
quality. This single sweep reproduces, within one dataset and one held-fixed artifact
detector, the same qualitative transition otherwise only visible by comparing across
GenBuster/SD~1.4/BigGAN or across the in-domain vs.\ cross-family-deployment
experiments (Section~VI.B).
%DIF >  NOTE: Table sourced from routing/results/e6_limited_data.csv, generated by
%DIF >  experiments/e6_limited_data.py (fixed this session to use FFT instead of
%DIF >  CNNSpot as the artifact detector -- see claude_context_notes/paper_revision_edits.txt
%DIF >  Section H/e6 note). GenBuster/BigGAN summary claims above are also drawn
%DIF >  from that CSV (dataset=="GenBuster"/"BigGAN" rows) but not tabulated in
%DIF >  full here for space; consider a supplementary table or figure if space
%DIF >  allows.
}

\DIFaddend \subsubsection{Cross-Dataset Bandit Transfer}

Table~\ref{tab:transfer} reports cross-dataset bandit transfer accuracy\DIFaddbegin \DIFadd{, now trained on
the same FFT+SigLIP2+Qwen features as the disagreement rule it is compared against,
correcting the CNNSpot+CLIP bandit vs.\
FFT+SigLIP2 disagreement-rule mismatch in the previous version of this comparison}\DIFaddend .
In-domain diagonal entries range from \DIFdelbegin \DIFdel{0.782 on GenBuster to 0.923 on BigGAN}\DIFdelend \DIFaddbegin \DIFadd{0.867 on SD~1.4 to 0.998 on BigGAN, matching the
disagreement rule almost exactly on every dataset}\DIFaddend . The bandits transfer across datasets
with off-diagonal accuracy ranging from \DIFdelbegin \DIFdel{0.650 to 0.896. BigGAN is }\DIFdelend \DIFaddbegin \DIFadd{0.760 to 0.998; BigGAN remains }\DIFaddend the strongest
transfer target \DIFdelbegin \DIFdel{, reflecting its cleaner verifier signal. }\DIFdelend \DIFaddbegin \DIFadd{and }\DIFaddend SD~1.4 \DIFdelbegin \DIFdel{is the weakest ,
where both out-of-domain bandits score below 0.75}\DIFdelend \DIFaddbegin \DIFadd{the weakest source}\DIFaddend , consistent with the noisy Qwen signal on
that benchmark.

The disagreement rule (FFT+SigLIP2, Qwen verifier\DIFaddbegin \DIFadd{; bottom row of
Table~\ref{tab:transfer}}\DIFaddend ) achieves 0.905, \DIFdelbegin \DIFdel{0.867, and }\DIFdelend 0.998\DIFdelbegin \DIFdel{across the three datasets (Table~\ref{tab:disagree_transfer}), exceeding }\DIFdelend \DIFaddbegin \DIFadd{, and 0.867 on GenBuster, BigGAN, and
SD~1.4 respectively -- at or above }\DIFaddend every cell of the \DIFaddbegin \DIFadd{fair }\DIFaddend bandit transfer matrix \DIFdelbegin \DIFdel{. The bandit reported here was trained and evaluated against
the }\DIFdelend \DIFaddbegin \DIFadd{(max
off-diagonal 0.897, still below GenBuster's 0.905 disagreement-rule accuracy), despite
having no trainable parameters and requiring no per-dataset retraining, unlike the
bandit. With a fair backbone, the in-domain bandit diagonal now essentially reproduces
the disagreement rule exactly on two of three datasets (SD~1.4: 0.867 vs.\ 0.867;
BigGAN: 0.998 vs.\ 0.998), a stronger version of claim C4 than under the previous
}\DIFaddend CNNSpot+CLIP \DIFdelbegin \DIFdel{configuration. Ordering is expected to hold with an FFT-based
bandit, as disagreement captures the dominant routing signal regardless of which
artifact detector is used. The disagreement rule does not require per-dataset retraining}\DIFdelend \DIFaddbegin \DIFadd{comparison, where the bandit trailed disagreement by 11.9 points on
SD~1.4}\DIFaddend .


\begin{table}[h]
\centering
\caption{Cross-dataset bandit transfer accuracy \DIFaddbeginFL \DIFaddFL{using FFT+SigLIP2+Qwen features (same
backbone as the disagreement rule)}\DIFaddendFL . Diagonal entries
are in-domain results. Off-diagonal entries apply a bandit trained on one dataset
directly to another without retraining. \DIFaddbeginFL \DIFaddFL{GenBuster's row/column uses a bandit fit on an
internal 70/30 split of its own 400-sample validation set, since Qwen predictions do not
cover GenBuster's training-split frames (see footnote). Bottom row is the disagreement
rule (no training, no per-dataset transfer -- included here rather than as a separate
table for direct comparison).}\DIFaddendFL }
\label{tab:transfer}
\begin{tabular}{lccc}
\toprule
Train $\rightarrow$ Test & GenBuster & BigGAN & SD14 \\
\midrule
GenBuster\DIFaddbeginFL \DIFaddFL{$^\dagger$ }\DIFaddendFL & \textbf{\DIFdelbeginFL \DIFdelFL{0.782}\DIFdelendFL \DIFaddbeginFL \DIFaddFL{0.900}\DIFaddendFL } & \DIFdelbeginFL \DIFdelFL{0.842 }\DIFdelendFL \DIFaddbeginFL \DIFaddFL{0.998 }\DIFaddendFL & \DIFdelbeginFL \DIFdelFL{0.650 }\DIFdelendFL \DIFaddbeginFL \DIFaddFL{0.867 }\DIFaddendFL \\
BigGAN    & \DIFdelbeginFL \DIFdelFL{0.770 }\DIFdelendFL \DIFaddbeginFL \DIFaddFL{0.843 }\DIFaddendFL & \textbf{\DIFdelbeginFL \DIFdelFL{0.923}\DIFdelendFL \DIFaddbeginFL \DIFaddFL{0.998}\DIFaddendFL } & \DIFdelbeginFL \DIFdelFL{0.733 }\DIFdelendFL \DIFaddbeginFL \DIFaddFL{0.760 }\DIFaddendFL \\
SD14      & \DIFdelbeginFL \DIFdelFL{0.730 }%DIFDELCMD < & %%%
\DIFdelFL{0.896 }\DIFdelendFL \DIFaddbeginFL \DIFaddFL{0.897 }\DIFaddendFL & \DIFdelbeginFL \textbf{\DIFdelFL{0.792}} %DIFAUXCMD
%DIFDELCMD < \\
%DIFDELCMD < \bottomrule
%DIFDELCMD < \end{tabular}
%DIFDELCMD < \end{table}
%DIFDELCMD < 

%DIFDELCMD < \begin{table}[h]
%DIFDELCMD < \centering
%DIFDELCMD < %%%
%DIFDELCMD < \caption{%
{%DIFAUXCMD
\DIFdelFL{Cross-dataset disagreement rule accuracy (FFT+SigLIP2, Qwen verifier). The
disagreement threshold has no trainable parameters and requires no dataset-specific
retraining.}}
%DIFAUXCMD
%DIFDELCMD < \label{tab:disagree_transfer}
%DIFDELCMD < \begin{tabular}{lc}
%DIFDELCMD < \toprule
%DIFDELCMD < %%%
\DIFdelFL{Dataset }\DIFdelendFL \DIFaddbeginFL \DIFaddFL{0.998 }\DIFaddendFL & \DIFdelbeginFL \DIFdelFL{Disagreement threshold accuracy }\DIFdelendFL \DIFaddbeginFL \textbf{\DIFaddFL{0.867}} \DIFaddendFL \\
\midrule
\DIFdelbeginFL \DIFdelFL{GenBuster }\DIFdelendFL \DIFaddbeginFL \DIFaddFL{Disagreement rule (no training) }\DIFaddendFL & 0.905 \DIFdelbeginFL %DIFDELCMD < \\
%DIFDELCMD < %%%
\DIFdelFL{SD~1.4    }%DIFDELCMD < & %%%
\DIFdelFL{0.867 }%DIFDELCMD < \\
%DIFDELCMD < %%%
\DIFdelFL{BigGAN    }\DIFdelendFL & 0.998 \DIFaddbeginFL & \DIFaddFL{0.867 }\DIFaddendFL \\
\bottomrule
\end{tabular}
\end{table}
\DIFaddbegin {\footnotesize \DIFadd{$^\dagger$GenBuster's Qwen predictions only cover its held-out
validation split (400 samples), not the training-split frames the standard protocol
uses to fit the bandit (as for SD14/BigGAN); this row/column instead uses an internal
70/30 split of the validation set and is not directly comparable to the other two rows
without a Qwen rerun on GenBuster's training split.}}
%DIF >  NOTE: Table regenerated by routing/bandit_transfer_fft_siglip2.py, source
%DIF >  routing/bandit_transfer_fft_siglip2.csv (this session). Former separate
%DIF >  tab:disagree_transfer table merged into this one as a bottom row this
%DIF >  session, addressing reviewer point (7). See
%DIF >  claude_context_notes/paper_revision_edits.txt Section D.
\DIFaddend 


\subsubsection{Sensitivity to CLIP Probe Quality}

The disagreement signal assumes both detectors provide meaningful predictions. When one detector degrades to near-random output, disagreement becomes uninformative.Table~\ref{tab:clip_transfer} shows that CLIP probes trained on one generator family do not generalize to others. A GenBuster-trained probe evaluated on SD~1.4 achieves 0.597
accuracy, near chance. Off-diagonal accuracies range from 0.505 to 0.598, showing that
cross-family probe transfer is unreliable regardless of source domain. Deploying the
cascade on a new generator family requires training a semantic probe on in-domain data.
The disagreement mechanism itself introduces no additional parameters, but depends on
both detectors providing signal above chance.


\begin{table}[h]
\centering
\caption{Cross-dataset transfer accuracy of the CLIP logistic probe. Diagonal entries
are in-domain. Off-diagonal entries apply the probe to another dataset without
retraining}

\label{tab:clip_transfer}
\begin{tabular}{lccc}
\toprule
Train $\rightarrow$ Test & GenBuster & SD~1.4 & BigGAN \\
\midrule
GenBuster & \textbf{0.965} & 0.597 & 0.511 \\
SD~1.4    & 0.598 & \textbf{0.913} & 0.551 \\
BigGAN    & 0.549 & 0.505 & \textbf{0.994} \\
\bottomrule
\end{tabular}
\end{table}

\DIFaddbegin \subsubsection{\DIFadd{Cheap-Trigger Ablation}}

\DIFadd{FFT contributes nothing on agreement cases by construction: the cascade falls back to
the semantic probe's own prediction there. Genuine value from FFT therefore requires
that its disagreements with the semantic probe pick out cases where escalation helps,
not merely some fixed-size subset of inputs. We test this directly on the three
in-domain benchmarks: at the same escalation rate as the real disagreement rule, we
instead escalate a random subset of the same size to Qwen, holding cost fixed, and
compare cascade accuracy over 2}{\DIFadd{,}}\DIFadd{000 random draws (independent recomputation of the
FFT+SigLIP2+Qwen cascade; point estimates differ from Table~\ref{tab:main_results} by
under a point due to probe-loading order, not a different pipeline).
}

\DIFadd{On SD~1.4, the disagreement-routed cascade (0.867) falls above the random-trigger
mean (0.839, 95\% interval $[0.817, 0.860]$): disagreement identifies a subset where
the semantic probe is comparatively weaker (82.6\% accuracy vs.\ 92.7\% overall), so
escalating those specific cases costs less accuracy than escalating a typical random
subset of the same size would. BigGAN shows the same pattern at its much lower
escalation rate (0.998 vs.\ 0.994 $[0.990, 0.996]$). On GenBuster, the disagreement-routed
cascade (0.897) falls inside the random-trigger interval (0.904, $[0.880, 0.925]$):
semantic-probe accuracy on the disagreement subset (95.5\%) is close to its overall
accuracy (98.5\%), so FFT's disagreements with SigLIP2 on GenBuster do not identify a
comparatively unreliable subset there, and escalating them is statistically
indistinguishable from escalating an equally-sized random sample. FFT's marginal value
over a same-cost arbitrary trigger holds on SD~1.4 and BigGAN but is not established on
GenBuster, consistent with GenBuster already being the benchmark where escalation is
least useful (Case 2, Section~V.B): the same failure to identify a genuinely harder
subset explains both results.
%DIF >  NOTE: routing/experiments/e11_cheap_trigger_ablation.py,
%DIF >  routing/results/e11_cheap_trigger_ablation.csv (this session). Addresses
%DIF >  Reviewer 2's request to quantify FFT's marginal value against a cheaper trigger.
}

\DIFaddend \subsubsection{Cross-Family Deployment}
\label{sec:cross_family_deployment}

Cross-family deployment evaluates whether the cascade provides benefit when the semantic
probe is applied out of domain. A probe trained on dataset A is applied to dataset B
without retraining, then compared against the stale probe alone, the \DIFaddbegin \DIFadd{FFT-only detector
(trained in-domain on the test dataset, as in Table~\ref{tab:main_results}),
and the }\DIFaddend cascade with Qwen as verifier \DIFdelbegin \DIFdel{, and an in-domain probe as an upper bound }\DIFdelend (Table~\ref{tab:cross_family_deployment}).

\DIFaddbegin \DIFadd{Adding the FFT-only reference column changes the interpretation of this experiment. }\DIFaddend The
cascade improves over the stale cross-family \DIFdelbegin \DIFdel{probe in four }\DIFdelend \DIFaddbegin \DIFadd{SigLIP2 probe in all six pairings (mean
$+$17.5 points) -- but this recovery is almost entirely attributable to the strength of
the FFT detector itself, not to the verifier or the disagreement mechanism: the cascade
beats FFT-only in only one }\DIFaddend of six pairings \DIFdelbegin \DIFdel{, with the largest gains when the
probe is near chance. Disagreement remains informative in these
cases and Qwen maintains above-chance }\DIFdelend \DIFaddbegin \DIFadd{(mean $-$10.2 points relative to FFT-only).
A paired bootstrap (10,000 resamples per cell) confirms all six $\Delta_{\text{stale}}$
values are statistically distinguishable from zero (narrowest 95\% CI: $+$0.090
$[+0.056, +0.123]$, BigGAN$\to$SD~1.4). $\Delta_{\text{FFT}}$ is a different picture:
four of six cells are significantly negative, but the two involving GenBuster as the
test set (SD~1.4$\to$GenBuster: $-$0.005 $[-0.048, +0.040]$; BigGAN$\to$GenBuster:
$+$0.003 $[-0.037, +0.043]$) have CIs that straddle zero, so the cascade is not
reliably distinguishable from FFT-only there either. The ``beats FFT-only in one of
six'' point estimate is itself noise-level rather than a second, weaker positive result.
A badly-transferred semantic probe creates a large, spurious ``cascade gain'' over
itself, while the informative comparison -- cascade vs.\ the stronger of the two cheap
detectors alone -- shows escalation providing no net benefit and typically hurting. This
is the sharpest empirical illustration of }\textbf{\DIFadd{Case 2}} \DIFadd{in the paper (Section~III.D):
in five of six pairings $A_{2}$ (stale SigLIP2) is so degraded that disagreement mostly
flags cases the FFT detector already gets right, and escalating those to Qwen trades
correct FFT predictions for Qwen's mixed }\DIFaddend accuracy on the routed subset\DIFdelbegin \DIFdel{, recovering up to 19
percentage points. The only negative case occurs when the verifier provides no useful
signal on the routed subset. An in-domain semantic probe remains the strongest choice
when available. When only a cross-family probe is available, the cascade provides a mean
gain of $+$6.3 points with no retraining}\DIFdelend \DIFaddbegin \DIFadd{. Practical
implication: under cross-family deployment shift, if a strong, architecture-general
artifact detector such as FFT is available, deploying it alone outperforms escalating a
cross-family-transferred semantic probe's disagreements to this verifier. This is a
different regime from the SD~1.4-to-SDXL/FLUX deployment shift in
Table~\ref{tab:e8_results}, which uses GPT-5.5 as verifier and shows genuine positive
cascade gain (Case 3) -- the two experiments should not be conflated: here the semantic
probe is transferred across generator }\emph{\DIFadd{families}} \DIFadd{(GenBuster/SD~1.4/BigGAN, trained
on data from one family and applied to another) with Qwen as verifier, whereas
Table~\ref{tab:e8_results} holds the family fixed (SD~1.4-trained probes) and shifts to
newer models from the }\emph{\DIFadd{same}} \DIFadd{diffusion lineage (SDXL/FLUX) with GPT-5.5 as verifier}\DIFaddend .

\DIFdelbegin %DIFDELCMD < \begin{table}[t]
%DIFDELCMD < %%%
\DIFdelendFL \DIFaddbeginFL \begin{table*}[t]
\DIFaddendFL \centering
\caption{Cross-family deployment results. Cascade uses the cross-family \DIFdelbeginFL \DIFdelFL{CLIP }\DIFdelendFL \DIFaddbeginFL \DIFaddFL{SigLIP2 }\DIFaddendFL probe
with Qwen escalation on disagreement cases \DIFaddbeginFL \DIFaddFL{(FFT is trained in-domain on the test
dataset, not transferred, matching Table~\ref{tab:main_results}).
$\Delta_{\text{stale}}$ = cascade $-$ stale SigLIP2; $\Delta_{\text{FFT}}$ = cascade $-$
FFT-only. In-domain SigLIP2 (ceiling): GenBuster 0.985, SD~1.4 0.927, BigGAN 0.996.
$\dagger$: 95\% bootstrap CI (10,000 resamples) includes zero; exact intervals in
Section~V.D.6}\DIFaddendFL .}
\label{tab:cross_family_deployment}
\footnotesize
\setlength{\tabcolsep}{4pt}
\DIFdelbeginFL %DIFDELCMD < \begin{tabular}{llcccc}
%DIFDELCMD < %%%
\DIFdelendFL \DIFaddbeginFL \begin{tabular}{llccccc}
\DIFaddendFL \toprule
Train & Test & \DIFdelbeginFL \DIFdelFL{Stale }\DIFdelendFL \DIFaddbeginFL \DIFaddFL{FFT }& \DIFaddFL{Stale SigLIP2 }\DIFaddendFL & Casc. & \DIFdelbeginFL \DIFdelFL{$\Delta$ }\DIFdelendFL \DIFaddbeginFL \DIFaddFL{$\Delta_{\text{stale}}$ }\DIFaddendFL & \DIFdelbeginFL \DIFdelFL{In-dom }\DIFdelendFL \DIFaddbeginFL \DIFaddFL{$\Delta_{\text{FFT}}$ }\DIFaddendFL \\
\midrule
GenBuster & SD~1.4    & \DIFdelbeginFL \DIFdelFL{0.546 }\DIFdelendFL \DIFaddbeginFL \DIFaddFL{0.748 }\DIFaddendFL & \DIFdelbeginFL \DIFdelFL{0.571 }\DIFdelendFL \DIFaddbeginFL \DIFaddFL{0.552 }& \DIFaddFL{0.654 }\DIFaddendFL & $+$\DIFdelbeginFL \DIFdelFL{0.025 }\DIFdelendFL \DIFaddbeginFL \DIFaddFL{0.102 }\DIFaddendFL & \DIFdelbeginFL \DIFdelFL{0.900 }\DIFdelendFL \DIFaddbeginFL \DIFaddFL{$-$0.094 }\DIFaddendFL \\
GenBuster & BigGAN    & \DIFdelbeginFL \DIFdelFL{0.523 }\DIFdelendFL \DIFaddbeginFL \DIFaddFL{0.998 }\DIFaddendFL & \DIFdelbeginFL \DIFdelFL{0.713 }\DIFdelendFL \DIFaddbeginFL \DIFaddFL{0.554 }& \DIFaddFL{0.787 }\DIFaddendFL & $+$\DIFdelbeginFL \DIFdelFL{0.190 }\DIFdelendFL \DIFaddbeginFL \DIFaddFL{0.233 }\DIFaddendFL & \DIFdelbeginFL \DIFdelFL{0.990 }\DIFdelendFL \DIFaddbeginFL \DIFaddFL{$-$0.210 }\DIFaddendFL \\
SD~1.4    & GenBuster & \DIFdelbeginFL \DIFdelFL{0.613 }\DIFdelendFL \DIFaddbeginFL \DIFaddFL{0.735 }\DIFaddendFL & \DIFdelbeginFL \DIFdelFL{0.562 }\DIFdelendFL \DIFaddbeginFL \DIFaddFL{0.547 }\DIFaddendFL & \DIFdelbeginFL \DIFdelFL{$-$0.050 }\DIFdelendFL \DIFaddbeginFL \DIFaddFL{0.730 }\DIFaddendFL & \DIFdelbeginFL \DIFdelFL{0.910 }\DIFdelendFL \DIFaddbeginFL \DIFaddFL{$+$0.182 }& \DIFaddFL{$-$0.005$^\dagger$ }\DIFaddendFL \\
SD~1.4    & BigGAN    & \DIFdelbeginFL \DIFdelFL{0.535 }\DIFdelendFL \DIFaddbeginFL \DIFaddFL{0.998 }\DIFaddendFL & \DIFdelbeginFL \DIFdelFL{0.721 }\DIFdelendFL \DIFaddbeginFL \DIFaddFL{0.613 }& \DIFaddFL{0.833 }\DIFaddendFL & $+$\DIFdelbeginFL \DIFdelFL{0.185 }\DIFdelendFL \DIFaddbeginFL \DIFaddFL{0.221 }\DIFaddendFL & \DIFdelbeginFL \DIFdelFL{0.990 }\DIFdelendFL \DIFaddbeginFL \DIFaddFL{$-$0.165 }\DIFaddendFL \\
BigGAN    & GenBuster & \DIFdelbeginFL \DIFdelFL{0.547 }\DIFdelendFL \DIFaddbeginFL \DIFaddFL{0.735 }\DIFaddendFL & \DIFdelbeginFL \DIFdelFL{0.547 }\DIFdelendFL \DIFaddbeginFL \DIFaddFL{0.517 }& \DIFaddFL{0.738 }\DIFaddendFL & $+$\DIFdelbeginFL \DIFdelFL{0.000 }\DIFdelendFL \DIFaddbeginFL \DIFaddFL{0.220 }\DIFaddendFL & \DIFdelbeginFL \DIFdelFL{0.910 }\DIFdelendFL \DIFaddbeginFL \DIFaddFL{$+$0.003$^\dagger$ }\DIFaddendFL \\
BigGAN    & SD~1.4    & \DIFdelbeginFL \DIFdelFL{0.508 }\DIFdelendFL \DIFaddbeginFL \DIFaddFL{0.748 }\DIFaddendFL & \DIFdelbeginFL \DIFdelFL{0.537 }\DIFdelendFL \DIFaddbeginFL \DIFaddFL{0.515 }& \DIFaddFL{0.604 }\DIFaddendFL & $+$\DIFdelbeginFL \DIFdelFL{0.029 }\DIFdelendFL \DIFaddbeginFL \DIFaddFL{0.090 }\DIFaddendFL & \DIFdelbeginFL \DIFdelFL{0.900 }\DIFdelendFL \DIFaddbeginFL \DIFaddFL{$-$0.144 }\DIFaddendFL \\
\midrule
\DIFdelbeginFL %DIFDELCMD < \multicolumn{2}{l}{Mean $\Delta$} %%%
\DIFdelendFL \DIFaddbeginFL \multicolumn{2}{l}{Mean} & \DIFaddendFL & & & $+$\DIFdelbeginFL \DIFdelFL{0.063 }\DIFdelendFL \DIFaddbeginFL \DIFaddFL{0.175 }\DIFaddendFL & \DIFaddbeginFL \DIFaddFL{$-$0.102 }\DIFaddendFL \\
\bottomrule
\end{tabular}
\DIFdelbeginFL %DIFDELCMD < \end{table}
%DIFDELCMD < %%%
\DIFdelend \DIFaddbegin \end{table*}
%DIF >  NOTE: Table regenerated with FFT+SigLIP2 (was CNNSpot+CLIP) by
%DIF >  routing/cross_family_deployment.py, source routing/cross_family_deployment.csv
%DIF >  (this session). This is the most consequential finding from this revision
%DIF >  pass -- see claude_context_notes/paper_revision_edits.txt Section E for
%DIF >  full context and required consistency edits elsewhere (Conclusion,
%DIF >  Practical Implications).
%DIF >  NOTE: dagger markers and CIs added this session from
%DIF >  routing/experiments/e10_bootstrap_deltas.py, routing/results/e10_cross_family_delta_cis.csv
%DIF >  (paired bootstrap, 10,000 resamples per cell). Addresses Reviewer 2's request for
%DIF >  uncertainty quantification on the delta columns.

\subsubsection{\DIFadd{A More Distant Shift: GAN Deployment}}
\label{sec:biggan_distant}

\DIFadd{The SDXL/FLUX deployment shift (Table~\ref{tab:e8_results}) and the cross-family
deployment shift above both stay within diffusion or video-diffusion generator
families. To test a genuinely more distant shift, we applied the SD~1.4-trained
(stale) FFT and SigLIP2 probes zero-shot to GenImage BigGAN, a GAN rather than a
diffusion model. Degradation is far more severe than under SDXL/FLUX: stale FFT
accuracy drops to 38.3\% and stale SigLIP2 to 62.1\% on the full 11}{\DIFadd{,}}\DIFadd{997-sample
pool, against 57.8--62.2\% / 81.5--92.0\% on SDXL/FLUX. This yields a 3}{\DIFadd{,}}\DIFadd{516-sample
disagreement set (29.3\%), from which we drew a stratified random sample of 1}{\DIFadd{,}}\DIFadd{500
(matched to the scale of the SDXL/FLUX disagreement subsets) and evaluated three
verifiers: GPT-5.5, Qwen2.5-VL-7B, and Qwen2.5-VL-72B.
}

\DIFadd{All three verifiers give a negative cascade gain over the stale semantic probe: GPT-5.5
$-$0.031 ($95\%$ CI $[-0.037, -0.026]$), Qwen2.5-VL-7B $-$0.064 ($[-0.071, -0.057]$), and
Qwen2.5-VL-72B $-$0.129 ($[-0.137, -0.122]$) -- even GPT-5.5, the only verifier that gave
consistent positive gain on SDXL/FLUX, fails here. The mechanism differs from the
cross-family Case~2 result above: on the disagreement subset specifically, stale
SigLIP2 accuracy is 90.8\% and stale FFT accuracy is 9.2\%, meaning disagreement here
mostly flags cases where the stale-transferred FFT probe is confused while the stale
semantic probe is still reliable, not cases that are genuinely hard for the semantic
probe. None of the three verifiers' disagreement-subset accuracy (GPT-5.5 80.1\%,
Qwen-7B 68.9\%, Qwen-72B 46.7\%) exceeds SigLIP2's 90.8\% there, so escalating to any of
them trades a likely-correct stale-probe answer for a less reliable one. Qwen-72B's
disagreement-subset accuracy is the lowest of the three and, notably, lower than the
smaller Qwen-7B model: a class breakdown shows Qwen-72B answering REAL on 97.5\% of
disagreement cases regardless of true label (98.6\% accuracy on the real class, 2.5\%
on the fake class), a severe real-answer bias on this generator family that a larger
model does not automatically avoid.
}

\DIFadd{This result directly validates the Verifier Validation Gate (Section~VI.I): a
calibration check on this disagreement subset would show every tested verifier's
accuracy below the stronger cheap detector's (stale SigLIP2's) accuracy on the same
subset, correctly recommending deployment of the stale semantic probe alone over
escalation, regardless of which of the three verifiers was under consideration. The
SDXL/FLUX Case~3 result does not generalize to this larger architectural shift: a
practitioner should not assume that escalation helps whenever both cheap detectors are
stale, only when a calibration check on the specific detector pair and verifier
confirms it.
%DIF >  NOTE: routing/run_gpt55_biggan_distant.py, routing/run_qwen72b_biggan.py,
%DIF >  routing/biggan_{gpt55,qwen7b,qwen72b}_stale_disagree_preds.csv (this session).
%DIF >  Addresses Reviewer 2's request for a more distant generator family
%DIF >  (GAN vs. diffusion, the reviewer's own suggested example) beyond SDXL/FLUX.
}\DIFaddend 

\begin{figure*}[t]
    \centering
    \includegraphics[width=\textwidth]{plot_fig_a_semantic_degradation.png}
    \caption{Semantic probe accuracy in-domain (solid) vs.\ deployment shift on SDXL and FLUX (hatched). In-domain accuracy reaches 86--99\%. On SDXL and FLUX it drops 5--17 points, with DINOv2 degrading most and SigLIP2 least. These stale probes define the deployment context where cascade routing provides gain.}

    \label{fig:semantic_degradation}
\end{figure*}


\section{Discussion}

This section interprets the empirical results presented in Section V. We first examine
why artifact-semantic disagreement separates easy from hard inputs\DIFaddbegin \DIFadd{, then map the
three K-stage regimes from Section~III.D onto the empirical results }\DIFaddend and assess whether
learned routing improves on \DIFdelbegin \DIFdel{this }\DIFdelend \DIFaddbegin \DIFadd{the disagreement }\DIFaddend signal, then characterize the role of the
semantic probe and VLM verifier under in-domain and deployment-shift conditions, and
finally discuss generator distribution shift, limitations, and practical deployment
guidelines.



\subsection{Artifact-Semantic Disagreement as a Routing Signal}

The FFT artifact detector achieves 99.7\%, 96.6\%, and 100\% accuracy on agreement
cases across the three benchmarks and drops to 4.9\%, 17.4\%, and 66.7\% on
disagreement cases. Agreement between the artifact and semantic detectors is a
near-perfect predictor of an easy sample, and disagreement predicts difficulty with
equal reliability. This pattern holds across fundamentally different generator
architectures: video-based diffusion generators (GenBuster), image diffusion models
(SD~1.4), and GAN-based generators (BigGAN).

The complementarity arises from the signal types each detector uses. The FFT probe
captures spectral regularities introduced during synthesis, including upsampling grids
and DCT coefficient patterns. Semantic probes encode high-level scene content learned
from natural image-text distributions. These signal types are largely orthogonal.
When both detectors agree, at least one type of signal is unambiguous and the sample
is easy to classify. When they disagree, neither is sufficient alone, indicating a
sample that requires higher-level reasoning. The disagreement rule routes only these
ambiguous cases to the VLM verifier without requiring a learned policy or calibrated
confidence scores.

BigGAN produces fewer than 1\% disagreement cases in-domain because GAN upsampling
artifacts are unambiguous in the frequency domain. The FFT probe resolves nearly all
BigGAN samples at the cheap stage. On GenBuster and SD~1.4, modern diffusion and
video generators produce 26--28\% disagreement rates, marking the regime where VLM
escalation recovers accuracy that neither cheap detector can provide alone.

\subsection{\DIFdelbegin \DIFdel{Limitations of Learned Routing}%DIFDELCMD < \MBLOCKRIGHTBRACE
%DIFDELCMD < %%%
\DIFdel{The retrained bandit router converges to the disagreement threshold on two of three
datasets. On GenBuster the bandit achieves 0.755 accuracy , identical to the
disagreement rule, }\DIFdelend \DIFaddbegin \DIFadd{Three Regimes in Practice}}
\label{sec:three_regimes}

\DIFadd{The K-stage formulation (Section~III.D) predicts three qualitative regimes, and each is
directly realized in the empirical results above. }\textbf{\DIFadd{Case 1}} \DIFadd{(both cheap stages
strong) is BigGAN: FFT alone reaches 99.8\% and disagreement collapses below 1\%, so the
cascade adds latency without changing accuracy (Table~\ref{tab:main_results}). }\textbf{\DIFadd{Case 2}}
\DIFadd{(asymmetric cheap stages) is the in-domain GenBuster and SD~1.4 result (Section~V.B) and,
more sharply, the cross-family deployment experiment
(Table~\ref{tab:cross_family_deployment}, Section~V.D.6): a badly cross-family-transferred
SigLIP2 probe creates a large apparent cascade gain over itself, but the cascade beats the
FFT detector alone in only one of six pairings, showing the gain is a detector-quality-gap
artifact rather than genuine routing value. }\textbf{\DIFadd{Case 3}} \DIFadd{(both cheap stages degraded) is
the SD~1.4-to-SDXL/FLUX shift with GPT-5.5 as verifier (Table~\ref{tab:e8_results}): both
the stale semantic probe and, to a lesser extent, FFT lose accuracy under this shift, and
escalation recovers up to 4.9 points with consistent positive gain across five of six
configurations (Figure~\ref{fig:cascade_gain}). A paired bootstrap on the four
configurations reconstructable from archived per-sample predictions (SigLIP2 }\DIFaddend and
\DIFdelbegin \DIFdel{on BigGAN it achieves 0.925, also identical}\DIFdelend \DIFaddbegin \DIFadd{DINOv2 on SDXL and FLUX) confirms three of the four gains are statistically
distinguishable from zero; the smallest, FLUX+SigLIP2 ($+$0.004), is not
($95\%$ CI $[-0.004, +0.012]$), so the weakest configuration in this Case~3 result
should be read as noise-level rather than a genuine (if small) recovery. Crucially, this Case~3 status is
verifier-dependent, not a property of the routing mechanism alone: substituting Qwen2.5-VL-7B
for GPT-5.5 as the verifier on the identical SDXL/FLUX disagreement subsets
(Table~\ref{tab:e8_results}) turns every one of the six configurations negative
($-$3.0 to $-$14.0 points), reverting the same degraded-cheap-detector setting to a
Case~2-like outcome. Both cheap detectors being weak is necessary but not sufficient for
Case 3; the verifier must also be reliably stronger than the stronger cheap detector on
the routed subset, exactly as Equation~}\eqref{eq:gain} \DIFadd{requires.
}

\DIFadd{Since GPT-5.5 is closed and proprietary, we additionally tested an open-weight strong
verifier, Qwen2.5-VL-72B-Instruct (4-bit quantized), on the same SDXL/FLUX SigLIP2
disagreement subsets used above. The result does not generalize uniformly across the
two datasets. On SDXL, Qwen-72B nearly reproduces GPT-5.5: disagreement-subset accuracy
of 89.8\% (vs.\ GPT-5.5's 90.0\%) and a cascade gain of $+$0.033 ($95\%$ CI
$[+0.024, +0.042]$), matching GPT-5.5's $+$0.034 gain on the same configuration within
noise. On FLUX, Qwen-72B behaves like the 7B model rather than GPT-5.5: disagreement-subset
accuracy of 58.5\% (vs.\ Qwen-7B's 58.5\% and GPT-5.5's 90.9\%) and a cascade gain of
$-$0.131 ($95\%$ CI $[-0.143, -0.118]$), reverting FLUX to a Case~2-like outcome under this
verifier. The open-weight reproducibility of the paper's Case~3 result is therefore
configuration-dependent: model scale alone (7B$\to$72B) recovers GPT-5.5-level performance
on SDXL but not on FLUX, so a practitioner substituting an open-weight verifier for GPT-5.5
should validate per-deployment rather than assume scale transfers the result.
%DIF >  NOTE: routing/run_qwen72b_sdxl_flux.py, routing/{sdxl,flux}_qwen72b_siglip2_disagree_preds.csv
%DIF >  (this session). Addresses Reviewer 2's request for an open-weight strong-verifier
%DIF >  alternative to GPT-5.5 for reproducibility independent of a proprietary API.
}

\DIFadd{The probe-degradation sweep (Table~\ref{tab:limited_data}, holding FFT fixed while the
semantic probe's training fraction shrinks from 100\% to 1\%) directly demonstrates the
Case~2-to-Case~3 transition as a continuous function of semantic-probe quality, rather
than relying on cross-dataset anecdote}\DIFaddend . On SD~1.4 \DIFdelbegin \DIFdel{the
bandit trails by
11.9 percentage points (0.623 vs.
\ 0.742) , where Qwen reward signal
is noisiest. These figures reflect the
CNNSpot+CLIP configuration.
}\DIFdelend \DIFaddbegin \DIFadd{with a DINOv2 probe, cascade gain over
semantic-only is $-$4.8 points at 100\% training data (Case 2) and crosses to $+$10.0
points at 1\% training data (Case 3), passing through near-zero around 5--10\% data; the
SigLIP2 backbone shows the same sign change between 1\% and 5\% data. On GenBuster, by
contrast, gain remains negative at every fraction tested, since even a 1\%-trained
semantic probe (16 labeled examples) already exceeds the FFT detector's accuracy on that
benchmark -- illustrating that the transition point depends on the FFT detector's own
strength on a given generator family, not on data quantity alone.
}

\DIFaddend The \DIFdelbegin \DIFdel{FFT-based
disagreement
rule achieves substantially higher accuracyon the same datasets
(Table~\ref{tab:crossdataset}), and the ordering is expected to hold}\DIFdelend \DIFaddbegin \textbf{\DIFadd{breakdown condition}} \DIFadd{(Section~III.D) -- disagreement degenerating when the two
cheap detectors share a correlated blind spot -- is realized directly in
Table~\ref{tab:escalation}: on GenBuster, the single case where FFT is wrong but the
semantic probe agrees with it (both detectors sharing the same error) is invisible to the
routing rule by construction, representing the fundamental limit of disagreement-based
routing at under 0.3\% of inputs.
}

\subsection{\DIFadd{Limitations of Learned Routing}}
\DIFadd{The bandit router, retrained on the same FFT+SigLIP2+Qwen features as the disagreement
rule it is compared against (Table~\ref{tab:transfer}), converges to the disagreement
threshold on all three datasets. On SD~1.4 the bandit achieves 0.867 accuracy, identical
to the disagreement rule; on BigGAN it achieves 0.998, also identical. On GenBuster the
bandit achieves 0.917 on an internal held-out split of its validation set (see
Table~\ref{tab:transfer} footnote), matching the disagreement rule's 0.917 on that same
split}\DIFaddend . A learned classifier with eleven detector features and explicit latency rewards
does not outperform the binary disagreement rule \DIFaddbegin \DIFadd{on any dataset}\DIFaddend , independently
validating that disagreement captures the dominant routing signal \DIFaddbegin \DIFadd{-- a stronger result
than under the previous CNNSpot+CLIP comparison, where the bandit trailed disagreement
by 11.9 points on SD~1.4 due to a mismatched backbone rather than a genuine policy gap}\DIFaddend .



\subsection{Role of the Semantic Probe}

The semantic probe (SigLIP2) achieves 98.5\%, 92.7\%, and 99.6\% accuracy at
12.9~ms cheap-stage latency across the three in-domain benchmarks. This is
Pareto-optimal in-domain where not invoking the VLM improved accuracy over the
semantic probe alone.

The FFT detector's role in the cascade is diagnostic, not predictive. On agreement
cases, the FFT and semantic probe produce identical predictions by construction. Any
routing policy that escalates only disagreement cases to the VLM reduces to outputting
the semantic probe's prediction on agreement cases unchanged. The FFT detector
 does not contribute to independent prediction on these samples. It serves as a cheap signal
whose agreement with the semantic probe identifies easy inputs, routing only the
residual disagreement cases to the VLM.

This clarifies when the cascade adds value. In-domain, the semantic probe is accurate
enough that VLM escalation does not improve accuracy and only adds latency. Under
deployment shift, when a stale probe loses accuracy on unseen generators, disagreement
rises and targeted VLM escalation recovers it. The SDXL and FLUX experiments confirm
this (Table~\ref{tab:e8_results}). Semantic probes do not transfer across generator
families (Table~\ref{tab:clip_transfer}), so production deployment on a new generator
family requires in-domain probe training. The disagreement rule itself needs no
retraining.

\subsection{VLM Domain Sensitivity}

By using the prompt mentioned in Section~IV.B, Qwen2.5-VL achieves a balanced accuracy of 50.8\% on GenBuster, 62.3\% on SD~1.4, and 73.4\% on BigGAN when applied uniformly to all inputs. This is a substantial improvement over text-decoded Qwen configurations we initially evaluated, which produced near-zero fake recall on static image datasets. We attribute this improvement to extracting the REAL/FAKE logit difference at generation time instead of parsing decoded output tokens. With standard greedy decoding, the model's text outputs were biased toward "REAL" even when its underlying logits distinguished between classes. This is a methodological observation that may generalize to other VLM-based classification tasks.

Despite this improvement, Qwen2.5-VL remains substantially weaker than the CLIP logistic probe on all three datasets. 
The disagreement-routed cascade extracts useful signal from Qwen by invoking it only
on inputs where the FFT and semantic probe disagree. These are precisely the inputs
where the cheap detectors are unreliable and any additional signal from the VLM is
valuable. On GenBuster, the cascade achieves 90.5\% accuracy even though uniform
Qwen achieves only 69.0\%.
 This demonstrates that selective escalation makes a mediocre verifier useful.

Although the uniform Qwen accuracy ranges from roughly 50-73\% across datasets, its accuracy on the routed disagreement subset is higher \DIFdelbegin \DIFdel{. In-domain, it reaches 58-75\% (Table~\ref{tab:verifier_comparison}), and }\DIFdelend \DIFaddbegin \DIFadd{on average than uniform escalation, though it varies widely by pairing: }\DIFaddend under cross-family deployment \DIFdelbegin \DIFdel{it achieves 60-68\% in 5 of 6 }\DIFdelend \DIFaddbegin \DIFadd{(FFT+SigLIP2 disagreement subset) it ranges 44.5--74.1\% across the six }\DIFaddend pairings (Table~\ref{tab:cross_family_deployment}). This \DIFdelbegin \DIFdel{shows that }\DIFdelend \DIFaddbegin \DIFadd{wide range -- and the fact that it is frequently below the FFT-only accuracy on the same pairing -- is why escalation does not reliably beat FFT-only in this setting (Section~V.D.6, Table~\ref{tab:cross_family_deployment}): }\DIFaddend the VLM provides useful signal on \DIFdelbegin \DIFdel{the samples it processes, even when overall accuracy is modest}\DIFdelend \DIFaddbegin \DIFadd{some routed subsets but not consistently enough to beat a strong cheap detector used alone}\DIFaddend .

\begin{figure}[h]
    \centering
    \includegraphics[width=\columnwidth]{plot_fig_b_cascade_gain_scatter.png}
    \caption{Cascade gain over the stale semantic probe vs.\ stale probe accuracy on SDXL and FLUX. Gain is negative only for FLUX + CLIP (91.9\% stale accuracy), where the probe remains strong enough that escalation cannot improve it. All five lower-accuracy configurations show positive gain, with the OLS trend confirming that cascade benefit scales inversely with stale probe quality.}

    \label{fig:cascade_gain}
\end{figure}

\subsection{Generator Distribution Shift}

CNNSpot accuracy is 55.0\%, 50.2\%, and 72.1\% across the three datasets. The
near-chance performance on GenBuster and SD~1.4 confirms that artifact-based detectors
trained on GANs do not generalize to modern diffusion and video
generators~\cite{Wang_2023_ICCV}. The FFT probe generalizes substantially better:
75.3\%, 74.8\%, and 99.8\%. Its spectral features capture regularities common across
synthesis architectures rather than GAN-specific upsampling patterns. The cascade
(FFT+SigLIP2, Qwen verifier) improves over FFT-only by 15.2 and 11.9 percentage
points on GenBuster and SD~1.4. On BigGAN, FFT-only already achieves 99.8\% and the
cascade matches it exactly.




\subsection{Limitations}

The evaluation sets are modest: 400 samples for GenBuster and 480 each for SD~1.4
and BigGAN. Bootstrap confidence intervals confirm statistical reliability at these
scales, but larger sets would strengthen generalization claims. The bandit router is
trained and evaluated within the same benchmark distribution, so its generalization
to unseen generator types remains untested. Qwen2.5-VL is applied zero-shot without
fine-tuning. However, a fine-tuned VLM may recover the domain sensitivity failures observed on
SD~1.4 and BigGAN. The study uses frame-level detection without temporal modeling,
and incorporating motion and frame consistency could improve performance on
video-specific generators. The semantic probe does not transfer across generator
families and deploying this cascade on a new family requires retraining the probe on
in-domain data. The disagreement rule itself has no trainable parameters and needs
no retraining.

\DIFaddbegin \DIFadd{The SDXL and FLUX deployment-shift evaluation (Section~IV.B) draws its real-image
half from ImageNet val/nature photographs, the same source used for the real
class of the SD~1.4 semantic probes' training data. Only the fake half of each
test set, the SDXL and FLUX generations, is genuinely unseen by the probes; the
real half is in-distribution by construction, so the reported deployment-shift
recovery is not fully independent of the probes' training distribution. A
class-conditional breakdown of the SigLIP2 configuration, representative of the
three semantic backbones in Table~\ref{tab:e8_results}, quantifies how much this
affects the result. On the real class, stale SigLIP2 accuracy is 95.8\% (SDXL) and
95.4\% (FLUX), consistent with in-distribution performance, and escalating
disagreement cases to GPT-5.5 lowers real-class accuracy by 4.0--4.3 points. On
the fake class, the genuinely out-of-distribution half, stale SigLIP2 accuracy
drops to 73.6\% (SDXL) and 86.5\% (FLUX), and the cascade raises it to 84.7\% and
91.2\% respectively. The net cascade gain in Table~\ref{tab:e8_results} comes
entirely from the fake class, at some cost to the real class, so real-image
alignment works against the reported gain rather than explaining it.
%DIF >  NOTE: class-conditional breakdown from routing/experiments/e9_realimage_independence.py,
%DIF >  routing/results/e9_class_breakdown.csv (this session). Addresses Reviewer 3's
%DIF >  AIGIBench real-image independence concern.
}

\DIFaddend \subsection{Practical Implications}

Disagreement-based routing provides a deployment-ready signal for adaptive detection
pipelines. It requires no calibration, no learned policy, and no retraining when the
artifact detector is already deployed. When a well-trained in-domain semantic probe
is available, deploying it without VLM escalation achieves the best cost-accuracy
tradeoff. When only a stale probe is available on unseen generators, routing disagreement cases to a frontier VLM recovers up to 4.9 percentage points
over a stale semantic probe while limiting escalation to 35--42\% of verifier
invocations (Table~\ref{tab:e8_results}, Table~\ref{tab:e8_efficiency}).
The cascade provides value \DIFaddbegin \DIFadd{only }\DIFaddend when the verifier is reliably stronger than the
\DIFaddbegin \emph{\DIFadd{stronger of the two cheap detectors}} \DIFadd{on the disagreement subset, not merely
stronger than a }\DIFaddend stale semantic probe\DIFdelbegin \DIFdel{on disagreement cases. When both are weak on the target domain, the stale probe alone }\DIFdelend \DIFaddbegin \DIFadd{. Under cross-family deployment specifically
(Table~\ref{tab:cross_family_deployment}), the FFT artifact detector alone
outperforms the full cascade in five of six pairings; a practitioner deploying on a
new generator family without an in-domain semantic probe should prefer the strong
cheap detector alone over escalating cross-family disagreements to Qwen. This
contrasts with the SD~1.4-to-SDXL/FLUX deployment-shift setting with GPT-5.5 as
verifier (Table~\ref{tab:e8_results}), where escalation does provide consistent
positive gain -- the practical guideline depends on both the verifier and the type of
shift, not on cascade routing being unconditionally beneficial under any distribution
shift. Under a still more distant shift (SD~1.4-trained probes applied zero-shot to
GenImage BigGAN, Section~V.D.7), even GPT-5.5 fails to provide positive gain: no
verifier's accuracy on the disagreement subset exceeds the stale semantic probe's own
90.8\% there, so a practitioner should not assume a strong verifier alone is sufficient
justification for escalation without first confirming it beats the stronger cheap
detector on the specific detector pair and target distribution in question.
}

\subsection{\DIFadd{A Practical Fail-Safe: the Verifier Validation Gate}}
\label{sec:validation_gate}

\DIFadd{A deployed system does not know in advance whether an incoming distribution is
in-domain, a benign generator update, or a hostile shift, so the guideline above --
prefer FFT-only under cross-family shift, prefer the GPT-5.5 cascade under SD~1.4-to-SDXL/FLUX
shift -- cannot be applied directly without first knowing which regime holds. A
natural first attempt is to use the disagreement rate itself as this signal, escalating
only when it exceeds the $\sim$26--28\% in-domain baseline (Section~V.C). This does not
work: the cross-family disagreement rate (Case~2, where escalation to Qwen hurts in five
of six pairings) ranges 38.5--50.4\%, and the SDXL/FLUX disagreement rate (Case~3, where
escalation to GPT-5.5 helps in five of six configurations) ranges 35.2--41.4\%
(Table~\ref{tab:e8_efficiency}) -- these ranges overlap substantially, so disagreement
rate alone cannot distinguish a regime where escalation helps from one where it hurts.
}

\DIFadd{We instead propose a }\textbf{\DIFadd{verifier validation gate}}\DIFadd{: before deploying escalation to a
verifier $V$ on the disagreement set of a new detector pair or a new target
distribution, evaluate $V$'s accuracy on a small labeled calibration sample drawn from
that disagreement set. Deploy the cascade only if $A_{V,\text{cal}} >
\max(A_1, A_2)$ measured on the same calibration sample; otherwise, fall back to
whichever cheap detector is individually stronger. This rule requires no knowledge of
whether the deployment }\DIFaddend is \DIFdelbegin \DIFdel{the preferred deployment }\DIFdelend \DIFaddbegin \DIFadd{in-domain or shifted -- it only requires a calibration sample
from wherever the system is currently running, which a practitioner can always collect
regardless of whether the distribution's identity is known.
}

\DIFadd{We validate this rule retrospectively against every configuration evaluated in this
paper where ground truth is available: the six cross-family pairings with Qwen
}\DIFaddend (Table~\ref{tab:cross_family_deployment})\DIFdelbegin \DIFdel{.
}\DIFdelend \DIFaddbegin \DIFadd{, the six SDXL/FLUX configurations under
each of GPT-5.5 (Table~\ref{tab:e8_results}) and Qwen (Table~\ref{tab:e8_results}), and
the three BigGAN distant-shift configurations under GPT-5.5, Qwen2.5-VL-7B, and
Qwen2.5-VL-72B (Section~V.D.7), 21 configurations total. Using the full
disagreement-subset accuracy as an oracle calibration sample, the gate's
escalate/fall-back decision matches the actual outcome (whether the cascade beat the
stronger cheap detector) in 20 of 21 configurations (95.2\%); the single miss, FLUX
with a SigLIP2 probe and GPT-5.5 as verifier, is a near-exact tie (90.88\% vs.\ 90.93\%,
a 0.04-point gap). The gate correctly identifies that Qwen should not be trusted as a
verifier in any of the twelve in-domain-family configurations tested (six cross-family,
six SDXL/FLUX), correctly identifies that GPT-5.5 should be trusted in five of six
SDXL/FLUX configurations, and correctly identifies that none of the three verifiers
should be trusted on the BigGAN distant shift, where even GPT-5.5's disagreement-subset
accuracy (80.1\%) falls well short of the stale semantic probe's (90.8\%) -- all without
ever being told which generator family, which verifier, or which shift type it was
evaluating. In deployment, the calibration sample would be smaller than the full
disagreement subset and introduce some estimation noise; we report the oracle version
here as an upper bound on the rule's achievable reliability, and leave
calibration-sample-size sensitivity to future work.
%DIF >  NOTE: Validated this session via routing/results/verifier_validation_gate_check.csv.
%DIF >  Directly addresses reviewer point 6 (practical fail-safe deployment guideline
%DIF >  given no prior knowledge of in-domain vs out-of-domain). The negative result
%DIF >  (disagreement rate alone does not separate Case 2 from Case 3) is reported
%DIF >  honestly rather than omitted, since it is what motivates the accuracy-based
%DIF >  gate instead of a simpler rate-based heuristic.
}\DIFaddend 




\section{Conclusion}

This work frames AI-generated media detection as an adaptive model selection problem. A two-stage cheap pipeline routes inputs through an FFT artifact detector and a semantic probe and escalates to a frontier VLM only when the two detectors disagree. Three research questions were evaluated with this design.

For RQ1, FFT and semantic probe agreement is a near-perfect predictor of detection ease. FFT accuracy is 99.7\%, 96.6\%, and 100\% on agreement cases and falls to 4.9\%, 17.4\%, and 66.7\% on disagreement cases. This pattern holds across video diffusion, image diffusion, and GAN generator families and forms the basis for the routing rule.

For RQ2, the disagreement rule does not require parameters or calibration. A bandit router trained with eleven detector features and explicit latency rewards\DIFaddbegin \DIFadd{, using the same FFT+SigLIP2+Qwen backbone as the disagreement rule, }\DIFaddend converges to the \DIFaddbegin \DIFadd{exact }\DIFaddend same routing decisions as the binary disagreement rule on \DIFdelbegin \DIFdel{two of three datasetsand trails by 11.9 percentage points on the third, where Qwen reward signal is noisiest}\DIFdelend \DIFaddbegin \DIFadd{all three datasets}\DIFaddend . This convergence independently validates that disagreement captures the dominant routing signal\DIFaddbegin \DIFadd{, and that more complex learned policies offer no measurable benefit over the parameter-free rule in this regime}\DIFaddend .

For RQ3, the cascade (FFT + SigLIP2, disagreement routing, Qwen verifier) achieves 90.5\%, 86.7\%, and 99.8\% accuracy on the three benchmarks. It escalates only 26--28\% of GenBuster and SD~1.4 inputs to the VLM and reduces mean latency by 62--70\% compared to always-escalate without sacrificing accuracy. Under deployment shift on SDXL and FLUX, the cascade recovers up to 4.9 percentage
points over a stale semantic probe while avoiding 58--65\% of verifier invocations
(Table~\ref{tab:e8_results}), directly satisfying RQ3's deployment-shift condition\DIFaddbegin \DIFadd{. This
positive result is specific to the SD~1.4-to-SDXL/FLUX shift with GPT-5.5 as verifier
(Case 3, Section~III.D); under cross-family shift across GenBuster/SD~1.4/BigGAN with
Qwen as verifier, escalation instead underperforms the FFT artifact detector alone in
five of six pairings (Table~\ref{tab:cross_family_deployment}), a Case 2 regime where a
badly-transferred semantic probe manufactures an illusory cascade gain. Testing an even
more distant shift, SD~1.4-trained probes applied zero-shot to GenImage BigGAN (a GAN,
architecturally unlike SD~1.4's diffusion process), no tested verifier -- including
GPT-5.5 -- gives positive cascade gain (Section~V.D.7), because disagreement there
mostly flags a stale, unreliable FFT probe rather than a genuinely hard case for the
semantic probe. RQ3's answer is
therefore conditional: disagreement-based escalation reduces cost while recovering
accuracy specifically when the verifier is reliably stronger than the }\emph{\DIFadd{stronger}} \DIFadd{of
the two cheap detectors on the routed subset, not whenever the semantic probe alone is
degraded, and not merely because both cheap detectors happen to be stale}\DIFaddend .


The effective deployment of heterogeneous detection pipelines depends on using cheap detectors to identify when expensive verification is needed rather than increasing individual detector complexity. Future work will explore temporal modeling for video detection, fine-tuned VLM verifiers for domain-specific forensic cues, and cross-family semantic probe training to remove the per-family retraining requirement.


\bibliographystyle{IEEEtran}
\bibliography{references}
\vskip -.5\baselineskip plus -1fil
\begin{IEEEbiography}[{\includegraphics[width=1in,height=1.25in,clip,keepaspectratio]{ansuman_nayak2.png}}]{Ansuman Nayak}
is a Software Engineer at Bank of America, where he focuses on building scalable automation tools for deployment pipelines, data workflows, and infrastructure operations. He earned his Bachelor's degree in Computer Science from Virginia Tech and is currently pursuing a Master's degree in Computer Science at Georgia Institute of Technology, specializing in Artificial Intelligence. His interests center on Agentic AI and LLM optimization, with a focus on workflow automation designed to operate efficiently under resource and cost constraints.
\end{IEEEbiography}
\vspace{-0.5\baselineskip}
\begin{IEEEbiography}[{\includegraphics[width=1in,height=1.25in,clip,keepaspectratio]{vijay.png}}]{Vijay K. Madisetti}
holds the position of Professor of Cybersecurity and Privacy (SCP) at Georgia Tech.
He earned his PhD in Electrical Engineering and Computer Sciences from the University of California at Berkeley
and is recognized as a Fellow of the IEEE. Additionally, he has been honored with the Terman Medal
by the American Society of Engineering Education (ASEE). Professor Madisetti has authored several widely
referenced textbooks on topics including cloud computing, data analytics, blockchain, and microservices.
\end{IEEEbiography}

\EOD

\end{document}
